% concatenated from cp_no_metrisable_quotient__3_.tex
\documentclass[a4paper,11pt]{amsart}
\usepackage{a4wide}
\usepackage[T1]{fontenc}
\usepackage[utf8]{inputenc}
\usepackage{lmodern}
\usepackage{amsmath,amssymb,mathtools,microtype,enumitem,url}
\usepackage{flafter,float}
\usepackage{tikz}
\usetikzlibrary{arrows.meta,positioning,calc}
\usepackage[colorlinks=true,linkcolor=blue,citecolor=blue,urlcolor=blue]{hyperref}
\setlength{\emergencystretch}{3em}
\theoremstyle{plain}
\newtheorem{theorem}{Theorem}[section]
\newtheorem{proposition}[theorem]{Proposition}
\newtheorem{lemma}[theorem]{Lemma}
\newtheorem{corollary}[theorem]{Corollary}
\newtheorem*{theoremA}{Theorem A}
\newtheorem*{theoremB}{Theorem B}
\theoremstyle{definition}
\newtheorem{definition}[theorem]{Definition}
\theoremstyle{remark}
\newtheorem{remark}[theorem]{Remark}
\newcommand{\R}{\mathbb R}
\newcommand{\N}{\mathbb N}
\newcommand{\supp}{\operatorname{supp}}
\newcommand{\Span}{\operatorname{span}}
\newcommand{\one}{\mathbf 1}
\newcommand{\eps}{\varepsilon}
\newcommand{\restr}[1]{{\restriction_{#1}}}
\numberwithin{equation}{section}
\title[A consistent failure of separable quotients]
{A consistent failure of separable quotients\protect\\for pointwise function spaces}
\author[T.~Kania]{Tomasz Kania}
\address[T.~Kania]{Mathematical Institute\\Czech Academy of Sciences\\\v Zitn\'{a} 25\\115 67 Praha 1\\Czech Republic and Institute of Mathematics and Computer Science\\Jagiellonian University\\{\L}ojasiewicza 6, 30-348 Krak\'{o}w, Poland}
\email{kania@math.cas.cz, tomasz.marcin.kania@gmail.com}
\thanks{RVO: 67985840.}
\author[J.~K\k{a}kol]{Jerzy K\k{a}kol}
\address[J.~K\k{a}kol]{Department of Functional Analysis\\Faculty of Mathematics and Computer Science\\Adam Mickiewicz University\\Uniwersytetu Pozna\'{n}skiego 4\\61-614 Pozna\'{n}, Poland}
\email{kakol@amu.edu.pl}
\date{}
\subjclass[2020]{Primary 54C35, 46A03; Secondary 03E35, 54D30, 28C15}
\keywords{Pointwise function space, separable quotient, metrisable quotient, Jensen diamond, inverse limit, finitely supported measure}
\begin{document}
\raggedbottom
\begin{abstract}
Assuming Jensen's diamond principle, we construct an infinite compact zero-dimensional space $K$ such that $C_p(K)$ has no infinite-dimensional Hausdorff separable linear quotient. The space $K$ is separable and crowded, has weight $\aleph_1$ and cardinality $2^{\aleph_1}$, and is an Efimov space. We construct $K$ as an inverse limit of compact metrisable spaces indexed by the countable ordinals. At each nontrivial successor step, the projection has two-point fibres over a chosen closed set and singleton fibres elsewhere; this changes the weak-star limit of a selected sequence of finitely supported measures. We also prove that, for compact $X$, existence of an infinite-dimensional separable quotient of $C_p(X)$ is equivalent to existence of an infinite-dimensional metrisable quotient.
\end{abstract}
\maketitle

\section{Introduction}

For a Tychonoff space $X$, let $C_p(X)$ denote the real vector space of continuous functions on $X$, endowed with pointwise convergence. K\k{a}kol and \'Sliwa \cite{KS2018} asked whether $C_p(K)$ has an infinite-dimensional separable quotient for every infinite compact space $K$. Banakh, K\k{a}kol and \'Sliwa formulated the corresponding metrisable-quotient question in \cite[Problem~2]{BKS2018}. We give a consistent negative answer to both questions. The construction excludes metrisable quotients; a compactness argument for finitely supported measures then excludes separable quotients.

These questions are pointwise counterparts of the Rosenthal--Lacey theorem: for infinite compact $K$, the Banach space $C(K)$ has a quotient isomorphic to $c_0$ or $\ell_2$; see \cite{Rosenthal,Lacey}. The pointwise problem has positive answers under several topological hypotheses. In particular, an infinite compact space containing a nontrivial convergent sequence or a copy of $\beta\N$ has an infinite-dimensional metrisable pointwise quotient \cite[Theorem~1 and Corollary~3]{BKS2018}. An \emph{Efimov space} is an infinite compact Hausdorff space containing neither of these two spaces. Thus any compact counterexample must be an Efimov space. For background, see Hart~\cite{Hart}.

Inverse-limit constructions of compact spaces under additional set-theoretic hypotheses have a substantial history. Fedorchuk's constructions \cite{Fedorchuk} are a principal antecedent of this approach, and Talagrand \cite{Talagrand} constructed, under CH, a compact space whose Banach space of continuous functions has the Grothendieck property: every weak-star convergent sequence in its dual is weakly convergent. We use successor projections with two-point fibres over a closed set, which may be uncountable, and singleton fibres elsewhere. The construction must preserve the measure convergence arranged at preceding stages; Lemma~\ref{lem:simultaneous-successor} provides this step.

Our main results are the following two theorems.

\begin{theoremA}
Assume $\diamondsuit_{\omega_1}$. There exists a separable crowded compact zero-dimensional space $K$ of weight $\aleph_1$ and cardinality $2^{\aleph_1}$ such that $C_p(K)$ has no infinite-dimensional Hausdorff separable linear quotient, and hence no infinite-dimensional Hausdorff metrisable linear quotient. In particular, $K$ is an Efimov space.
\end{theoremA}

The reduction from separable to metrisable quotients uses the following unconditional fact.

\begin{theoremB}
Let $X$ be compact Hausdorff. Then $C_p(X)$ has an infinite-dimensional Hausdorff separable linear quotient if and only if it has an infinite-dimensional Hausdorff metrisable linear quotient.
\end{theoremB}

All quotients in this paper carry their quotient topology, and all compact spaces are Hausdorff. The term \emph{crowded} means that the space has no isolated points. The principle $\diamondsuit_{\omega_1}$ is recalled in Section~\ref{sec:limits}. Since diamond holds in the constructible universe by Jensen's theorem \cite{Jensen}, the consistency of ZFC implies the consistency of ZFC together with the conclusion of Theorem~A; see also \cite[Theorem~1.2]{Rinot}. We do not obtain the conclusion from CH alone or from $\mathrm{MA}+\neg\mathrm{CH}$: the prediction used in the proof is stronger than CH, and CH is consistent with the failure of diamond \cite[Theorem~1.4]{Rinot}.

A related sufficient condition for a metrisable quotient is the
\emph{two-disjoint-copies property} (2DCP): there are nonempty compact
sets $(X_n)_{n<\omega}$ in $X$ such that each $X_n$ contains two
disjoint copies of $X_{n+1}$ \cite[Theorem~2]{BKS2018}.
The hereditarily weakly Koszmider compacta\footnote{This means that,
for every closed $L\subseteq X$, every bounded linear operator
$T:C(L)\to C(L)$ is a weak multiplier: $(Tf_n)(x_n)\to0$ whenever
$(f_n)\subseteq C(L)$ is bounded and pairwise disjoint,
$(x_n)\subseteq L$, and $f_n(x_n)=0$ for all $n$.}
of Barbeiro and Fajardo \cite{BarbeiroFajardo} are consistent Efimov
examples for which $C(X)$ is Grothendieck
\cite[Theorem~2.4]{Koszmider}. They fail 2DCP: interchanging two
disjoint infinite homeomorphic closed subspaces would give an operator
on their union which is not a weak multiplier.

In the companion paper \cite{KaniaKakol2DCP}, we use diamond to obtain
a Talagrand compactum $T$ which fails 2DCP while $C(T)$ remains
Grothendieck and its weak-star dual unit ball contains no copy of
$\beta\N$. That construction leaves the existence of an arbitrary
infinite-dimensional metrisable quotient of $C_p(T)$ open.
The present argument excludes such quotients by controlling the
weak-star closure of every countably infinite-dimensional subspace
of $L(K)$. We use direct diamond prediction as in the companion
construction; the additional clopen localisations will give the
nonzero finitely supported measures required for this argument.

We outline the proof of Theorem~A. A metrisable quotient of $C_p(K)$
would give a countably infinite-dimensional relatively weak-star closed
subspace $A$ of the finitely supported measures on $K$. The union of the
supports of its members is countable. Consequently, some projection onto
a countable stage is injective on this union and on $A$.
Diamond allows us to anticipate the projected subspace while constructing
that stage. We arrange that the weak-star closure of $A$ in $C(K)'$
contains a nonzero measure $\delta_y-\delta_x$ whose projection is zero.
Relative closedness would place this measure in $A$, contradicting
injectivity of the projection on $A$.

The main work is to preserve this obstruction through the remaining
stages. At the stage where a subspace is anticipated, we choose disjointly
supported signed measures whose positive and negative parts converge to
the same probability $\lambda$. We then choose a compact set $F$ with
$\lambda(F)>0$ and replace each $x\in F$ by two points $(x,0)$ and $(x,1)$.
For each clopen set $D$ in that stage, the new coordinates are chosen so
that a subsequence of the lifted signed measures converges, after
normalisation, to the difference of the two copies of the restriction
of $\lambda$ to $F\cap D$.

At subsequent stages, these limits are preserved uniformly over all
measures having the prescribed projection and norm. At countable limit
stages, a diagonal choice of subsequences preserves convergence. At the
final stage, each continuous function factors through a countable stage,
which suffices to retain the required weak-star closure relations.

The two copies of $F$ also determine a positive measure on pairs of
points. We extend this measure through the inverse system to a nonzero
Radon measure $\gamma$ on $K^2$. Its support consists of distinct pairs
with the same projection at the stage under consideration. The convergence
identities imply that every $f\in C(K)$ with $\mu(f)=0$ for all $\mu\in A$
satisfies
\[
\int_{K^2}|f(y)-f(x)|^2\,d\gamma(x,y)=0.
\]
By continuity, $f(x)=f(y)$ for every $(x,y)\in\supp\gamma$.
Thus any pair in this fixed nonempty support gives the required
$\delta_y-\delta_x$ in the weak-star closure of $A$.

Section~\ref{sec:prelim} gives the preliminary facts and proves Theorem~B.
Section~\ref{sec:limits} establishes the prediction and countable-limit
arguments. Section~\ref{sec:recursion} proves the successor lemma and
uses it to construct the inverse system and its measure limits.
Section~\ref{sec:coupling} converts the localised limits into differences
of point masses, which exclude quotients in Section~\ref{sec:finish}.
The remaining assertions of Theorem~A and the obstruction to using only
one-point splits are proved in Sections~\ref{sec:properties}
and~\ref{sec:obstruction}.

\section{Duality and preliminary facts}\label{sec:prelim}

We use $\N=\{0,1,2,\ldots\}$. For a compact space $X$, the norm on
$C(X)$ is $\|f\|_\infty=\sup_{x\in X}|f(x)|$. We identify $M(X)=C(X)'$
with the finite signed Radon measures and write
$\mu(f)=\int_X f\,d\mu$. A finite positive Radon measure $\lambda$ is
regular: for every Borel set $B\subseteq X$,
\[
\lambda(B)=\sup\{\lambda(F):F\subseteq B,\ F\text{ compact}\}
=\inf\{\lambda(V):B\subseteq V,\ V\text{ open}\}.
\]
The first equality is inner regularity and the second is outer regularity.
We write $P(X)$ for the Radon probability measures on $X$.

For $\mu\in M(X)$, its total variation is the positive measure
\[
|\mu|(B)=\sup\left\{\sum_{j=1}^r|\mu(B_j)|:
 B=\bigcup_{j=1}^r B_j,\ B_j\text{ pairwise disjoint Borel sets}\right\}.
\]
The variation norm is $\|\mu\|=|\mu|(X)$, equivalently
$\|\mu\|=\sup_{\|f\|_\infty\leqslant1}|\mu(f)|$.
The \emph{Jordan decomposition} is $\mu=\mu^+-\mu^-$, where
$\mu^+=(|\mu|+\mu)/2$ and $\mu^-=(|\mu|-\mu)/2$ are positive measures concentrated on disjoint Borel sets. In particular, if
$\mu=\sum_{j=1}^r a_j\delta_{x_j}$ with distinct $x_j$ and nonzero $a_j$,
then $\delta_x(f)=f(x)$,
\[
\mu^+=\sum_{a_j>0}a_j\delta_{x_j},\qquad
\mu^-=\sum_{a_j<0}(-a_j)\delta_{x_j},\qquad
\|\mu\|=\sum_{j=1}^r|a_j|.
\]
The space of these finitely supported measures is denoted by $L(X)$.

The \emph{support} of $\mu$ is
$\supp\mu=\{x\in X:|\mu|(V)>0\text{ for every open neighbourhood }V\ni x\}$.
We say that $\mu$ is \emph{concentrated on} a Borel set $E$ if
$|\mu|(X\setminus E)=0$. A positive Radon measure is \emph{nonatomic}
if it has no atoms; equivalently, it gives measure zero to every singleton.
A closed set is \emph{perfect} if it has no isolated points.
For a Borel set $E$, the restriction is the measure
$(\mu\restr{E})(B)=\mu(B\cap E)$. We denote its indicator function by
$\one_E$. For any subset $S\subseteq X$, put
$L(S)=\{\mu\in L(X):\supp\mu\subseteq S\}$; the set $S$ need not be closed.
The weak-star topology is $\sigma(M(X),C(X))$, denoted by $w^*$.
Thus $\mu_n\to\mu$ weak-star means $\mu_n(f)\to\mu(f)$ for every $f\in C(X)$.

If $\pi:X\to Y$ is continuous, its \emph{push-forward} on measures is
\[
(\pi_*\mu)(B)=\mu(\pi^{-1}[B]),\qquad
(\pi_*\mu)(f)=\mu(f\circ\pi)\quad(f\in C(Y)).
\]
It satisfies $\|\pi_*\mu\|\leqslant\|\mu\|$, since
$\|f\circ\pi\|_\infty\leqslant\|f\|_\infty$.
The same set formula defines $s_*\mu$ for a Borel map $s$;
when the spaces are compact metrisable, this is again a Radon measure.
The \emph{fibre} over $y$ is $\pi^{-1}\{y\}$. A \emph{Borel section} of
$\pi$ is a Borel map $s:Y\to X$ such that $\pi\circ s=\mathrm{id}_Y$.
We call $\zeta\in M(X)$ a \emph{norm-preserving lift} of $v\in M(Y)$
if $\pi_*\zeta=v$ and $\|\zeta\|=\|v\|$.

For a linear subspace $A\subseteq M(X)$ and a linear subspace $Z\subseteq C(X)$, define
\[
\begin{split}
A_\perp&=\{f\in C(X):\mu(f)=0\ (\mu\in A)\},\\
Z_M^\perp&=\{\mu\in M(X):\mu(f)=0\ (f\in Z)\},\qquad
Z_L^\perp=Z_M^\perp\cap L(X).
\end{split}
\]
The bipolar theorem for a dual pair gives
\begin{equation}\label{eq:bipolar}
\overline A^{\,w^*}=(A_\perp)_M^\perp,\qquad
\overline A^{\,w^*}\cap L(X)=(A_\perp)_L^\perp\quad(A\subseteq L(X)).
\end{equation}
Here the closure on both left-hand sides is in $M(X)$; intersecting with $L(X)$ gives the relative closure. These are standard duality facts; see \cite[Chapter~8]{Jarchow}. 

The first lemma expresses the metrisable-quotient question in terms of
finitely supported measures.

\begin{lemma}\label{lem:dual}
The space $C_p(X)$ has an infinite-dimensional metrisable Hausdorff quotient if and only if $L(X)$ contains a countably infinite-dimensional relatively weak-star closed linear subspace.
\end{lemma}
\begin{proof}
The topology of $C_p(X)$ is $\sigma(C(X),L(X))$; see \cite{Arhangelskii}.
For every closed linear subspace $Z\subseteq C_p(X)$, the quotient topology is
\[
\sigma(C(X)/Z,Z_L^\perp),
\]
by the standard quotient theorem for weak topologies \cite[Chapter~8]{Jarchow}.
Its continuous dual is therefore $Z_L^\perp$.

An infinite-dimensional metrisable quotient $Q$ is linearly homeomorphic
to a dense vector subspace of $\R^\N$. Indeed, a countable neighbourhood
base for its weak topology uses only countably many functionals. These
span $Q'$, since every continuous linear functional depends on finitely
many of the functionals defining a weak topology. Choose a basis
$(\lambda_n)_{n<\omega}$ of $Q'$. The map
$q\mapsto(\lambda_n(q))_{n<\omega}$ is a topological embedding with dense
image: each finite coordinate projection is surjective by linear
independence of the corresponding functionals. Conversely, every dense
vector subspace of $\R^\N$ is infinite dimensional and metrisable.
Thus $C_p(X)$ has such a quotient precisely when it admits a continuous
open linear surjection $T:C_p(X)\to D$ onto such a subspace.

The continuous dual $D'$ is the span of the coordinate evaluations,
which are linearly independent by density. Since $T$ is open and
surjective, its adjoint is injective and
\[
T'(D')=(\ker T)_L^\perp.
\]
This is a countably infinite-dimensional relatively weak-star closed
subspace of $L(X)$.

Conversely, let $A\subseteq L(X)$ have these properties and put
$Z=A_\perp$. Then $Z$ is closed in $C_p(X)$ and \eqref{eq:bipolar} gives
$Z_L^\perp=A$. The quotient topology displayed above is generated by a
countable basis of $A$, so it is metrisable. Its dual is infinite
dimensional, and hence so is the quotient.
\end{proof}

To apply the successor construction to a subspace of $L(X)$, we need
a sequence of disjointly supported signed measures whose positive and
negative parts have a common weak-star limit.

\begin{lemma}\label{lem:flat}
Let $X$ be compact metrisable and let $A$ be an infinite-dimensional linear subspace of $L(X)$. There exist $\sigma_n\in A$ and $\lambda\in P(X)$ such that the supports of the $\sigma_n$ are pairwise disjoint,
\[
\|\sigma_n\|=2,\qquad \sigma_n(\one)=0,\qquad
\sigma_n\longrightarrow0\quad\hbox{weak-star},
\]
and the Jordan parts $p_n=\sigma_n^+$ and $q_n=\sigma_n^-$ are probabilities converging weak-star to $\lambda$.
\end{lemma}
\begin{proof}
Choose a norm-dense sequence $f_0=\one,f_1,f_2,\ldots$ in $C(X)$. For $x\in X$, the coefficient map $a_x:L(X)\to\R$, $a_x(\mu)=\mu(\{x\})$, is an algebraic linear functional. At step $n$, intersect $A$ with the kernels of $\mu\mapsto\mu(f_j)$ for $j\leqslant n$, and with $\ker a_x$ for the finitely many points used in the earlier supports. The intersection has finite codimension in $A$, so it remains infinite dimensional. Choose a vector of norm $2$ in it. Its support is disjoint from all preceding supports.

For each $j$, the construction gives $\sigma_n(f_j)=0$ whenever $n\geqslant j$.
Given $f\in C(X)$ and $\eps>0$, choose $j$ with
$\|f-f_j\|_\infty<\eps/2$. For $n\geqslant j$,
\[
|\sigma_n(f)|=|\sigma_n(f-f_j)|
\leqslant\|\sigma_n\|\,\|f-f_j\|_\infty<\eps.
\]
Hence $\sigma_n\to0$ weak-star. Since each measure has total mass zero and norm $2$, both Jordan parts have mass one. The space $P(X)$ is weak-star compact and metrisable, so a subsequence of the positive parts converges to some $\lambda\in P(X)$. The differences tend to zero, hence the negative parts have the same limit on this subsequence. Reindex the subsequence to obtain all the stated properties. The measure $\lambda$ may have atoms; no subsequent argument requires it to be nonatomic.
\end{proof}

The next lemma controls all norm-preserving lifts of a finitely supported
measure. The concentration assertion will ensure uniqueness when the
fibres over its support are singletons.

\begin{lemma}\label{lem:lifts}
Let $\pi:X\to Y$ be a continuous surjection of compact Hausdorff spaces, and let $v\in L(Y)$ have norm $1$. Then
\[
\mathcal B(v,\pi)=\{\zeta\in M(X):\|\zeta\|\leqslant1,\ \pi_*\zeta=v\}
\]
is nonempty and weak-star compact. Every $\zeta\in\mathcal B(v,\pi)$
satisfies
\[
\|\zeta\|=1,\qquad |\zeta|\bigl(X\setminus\pi^{-1}(\supp v)\bigr)=0.
\]
For an arbitrary nonzero $v\in L(Y)$, the same conclusions hold with
$1$ replaced by $\|v\|$ in the definition of $\mathcal B(v,\pi)$ and in
the norm equality.
\end{lemma}
\begin{proof}
Choose one point above each atom of $v$ to obtain a norm-preserving lift. Compactness follows from Banach--Alaoglu and weak-star continuity of push-forward. For every $\zeta\in\mathcal B(v,\pi)$,
\[
1=\|v\|\leqslant\|\zeta\|\leqslant1,\qquad
|\pi_*\zeta|\leqslant\pi_*|\zeta|.
\]
The positive measures in the second inequality have equal total mass, so they are equal. Since $|v|$ is concentrated on $\supp v$, the equality $\pi_*|\zeta|=|v|$ gives $|\zeta|(X\setminus\pi^{-1}(\supp v))=0$.
\end{proof}

We recall the inverse-system terminology used in the construction; see
\cite[Sections~2.5 and 3.2]{Engelking} and
\cite[Section~20.5]{KKLPS}, together with the references therein.
An \emph{inverse system} $(K_\alpha,\pi_{\alpha,\beta})$, indexed by an
ordinal interval starting at $\alpha_0$, consists of compact spaces and continuous maps $\pi_{\alpha,\beta}:K_\beta\to K_\alpha$
for $\alpha\leqslant\beta$, with $\pi_{\alpha,\alpha}=\mathrm{id}$ and
$\pi_{\alpha,\gamma}=\pi_{\alpha,\beta}\circ\pi_{\beta,\gamma}$ whenever
$\alpha\leqslant\beta\leqslant\gamma$. Its inverse limit below a limit
ordinal $\rho>\alpha_0$ is the compact space
\[
\varprojlim_{\alpha_0\leqslant\alpha<\rho}K_\alpha
=\left\{(x_\alpha)\in\prod_{\alpha_0\leqslant\alpha<\rho}K_\alpha:
 \pi_{\alpha,\beta}(x_\beta)=x_\alpha\ (\alpha_0\leqslant\alpha\leqslant\beta<\rho)\right\}.
\]
The system is \emph{continuous} if, at each limit stage $\rho>\alpha_0$ in its
index set, the natural map from $K_\rho$ to this inverse limit is a
homeomorphism. In our systems the indices start at $\omega$, the maps
are surjective, and $K_\alpha\subseteq2^\alpha$, where
$2^\alpha=\{0,1\}^\alpha$ has the product topology. The maps are coordinate
restrictions, $\pi_{\alpha,\beta}(x)=x\restr{\alpha}$.

A continuous surjection is a \emph{simple extension} if exactly one fibre
has two points and all other fibres are singletons. In
\cite[Definition~6.1 and Theorem~6.10]{KSZ}, K\k{a}kol, Sobota and Zdomskyy
proved that every infinite limit $X$ of a continuous inverse system of
zero-dimensional compact spaces starting with a singleton and using
simple extensions at each successor has the \emph{finitely supported
Josefson--Nissenzweig property}. This means that there is a sequence
$(v_n)\subseteq L(X)$ such that $\|v_n\|=1$ and $v_n\to0$ weak-star;
such a sequence is called a finitely supported Josefson--Nissenzweig
sequence. Section~\ref{sec:obstruction} examines the restriction to
simple extensions and its relation to our construction.

The following factorisation fact allows us to check a continuous function
on the final space at a countable stage.

\begin{lemma}\label{lem:factor}
Let $(K_\alpha,\pi_{\alpha,\beta})_{\omega\leqslant\alpha\leqslant\beta\leqslant\omega_1}$ be a continuous inverse system of compact spaces with surjective bonding maps, where $K_\alpha\subseteq2^\alpha$ and the maps are coordinate projections. Then
\[
C(K_{\omega_1})=\bigcup_{\alpha<\omega_1,\ \alpha\geqslant\omega}
\pi_{\alpha,\omega_1}^*C(K_\alpha),\qquad
\pi_{\alpha,\omega_1}^*g=g\circ\pi_{\alpha,\omega_1}.
\]
\end{lemma}
\begin{proof}
Fix $f\in C(K_{\omega_1})$. For every positive integer $m$, continuity and compactness give a finite cover by relative clopen cylinders on which the oscillation of $f$ is less than $2^{-m}$. The union of the finite coordinate sets used in these covers is countable, and therefore contained in some $\omega\leqslant\alpha<\omega_1$. If $x,y\in K_{\omega_1}$ satisfy
$\pi_{\alpha,\omega_1}(x)=\pi_{\alpha,\omega_1}(y)$, they have the same
coordinates on every finite set used in the covers. For each $m$, any
cylinder in the $m$th cover containing $x$ therefore also contains $y$.
Thus $|f(x)-f(y)|<2^{-m}$ for every $m$, and $f(x)=f(y)$.
We may consequently define
$g(z)=f(x)$ for any $x\in\pi_{\alpha,\omega_1}^{-1}\{z\}$;
this is well defined and $f=g\circ\pi_{\alpha,\omega_1}$. The surjective projection is a quotient map from a compact space onto a Hausdorff space, so $g$ is continuous. Conversely, every displayed pullback is continuous.
\end{proof}


We now prove Theorem~B. The key point is that measures with uniformly
bounded norms and support sizes form weak-star compact sets.

\begin{proof}[Proof of Theorem B]
By the weak-topology description in Lemma~\ref{lem:dual}, a Hausdorff metrisable quotient embeds topologically into $\R^\N$ by evaluation against a countable spanning family of its dual. It is therefore second countable and separable.

Conversely, let $Q=C_p(X)/Z$ be an infinite-dimensional Hausdorff separable quotient and put $A=Z_L^\perp$. This subspace of $L(X)$ is relatively weak-star closed, and the topology of $Q$ is the weak topology determined by $A$. Choose a countable dense set in $Q$ and representatives $f_n\in C(X)$ of its members. These functions separate the elements of $A$: if $\mu\in A$ annihilates every $f_n$, its induced continuous functional on $Q$ vanishes on a dense set and hence is zero. In particular, the map
\[
J:A\longrightarrow\R^\N,\qquad J(\mu)=(\mu(f_n))_{n<\omega},
\]
is injective.

For positive integers $r,m$, put
\[
\Delta_{r,m}(X)=\{\mu\in L(X):|\supp\mu|\leqslant r,\ \|\mu\|\leqslant m\},
\qquad S_{r,m}=A\cap\Delta_{r,m}(X).
\]
The set $\Delta_{r,m}(X)$ is weak-star compact: it is the image of
\[
X^r\times\{(a_1,\ldots,a_r)\in\R^r:\textstyle\sum_{j=1}^r|a_j|\leqslant m\}
\]
under the continuous map $(x,a)\mapsto\sum_{j=1}^ra_j\delta_{x_j}$.
Since $A$ is relatively closed in $L(X)$ and $\Delta_{r,m}(X)\subseteq L(X)$, the set $S_{r,m}$ is weak-star compact. The restriction of $J$ to $S_{r,m}$ is a continuous injection into the Hausdorff space $\R^\N$, hence a homeomorphism onto its image. Thus every $S_{r,m}$ is weak-star metrisable.

Suppose first that some $S_{r,m}$ is not compact for the variation norm $\|\mu\|=|\mu|(X)$. The identity map from its compact weak-star topology to its norm topology cannot be continuous, since a continuous image of a compact space is compact. At a point of discontinuity, metrisability gives $\mu_n,\mu\in S_{r,m}$ and $\eps>0$ such that
\[
\mu_n\xrightarrow{w^*}\mu,\qquad \|\mu_n-\mu\|\geqslant\eps\quad(n<\omega).
\]
The finitely supported measures $u_n=(\mu_n-\mu)/\|\mu_n-\mu\|$ have norm one and are weak-star null: for every $f\in C(X)$,
$|u_n(f)|\leqslant\eps^{-1}|(\mu_n-\mu)(f)|\to0$.
By \cite[Theorem~1]{BKS2019}, $C_p(X)$ has a quotient isomorphic to $(c_0)_p$, the space $c_0$ with the topology inherited from $\R^\N$. This quotient is infinite dimensional, Hausdorff and metrisable.

Suppose instead that every $S_{r,m}$ is norm compact. Choose a countable norm-dense subset $D_{r,m}\subseteq S_{r,m}$ and let
\[
S=\bigcup_{r,m\geqslant1}\ \bigcup_{\nu\in D_{r,m}}\supp\nu.
\]
The set $S$ is countable. If $\mu\in S_{r,m}$ had a nonzero atom at a point $x\notin S$, then
$\|\mu-\nu\|\geqslant|\mu(\{x\})|>0$ for every $\nu\in D_{r,m}$, contradicting norm density. Hence every $\mu\in A=\bigcup_{r,m\geqslant1}S_{r,m}$ is supported in $S$. It follows that $A\subseteq L(S)$ has countable Hamel dimension. Lemma~\ref{lem:dual}, including its description of the quotient topology, now shows that the original quotient $Q$ is metrisable. This proves the converse implication.
\end{proof}

\section{Prediction and preservation at countable limits}\label{sec:limits}

We prepare two ingredients for the construction of $K$. First, diamond
will predict every countable sequence of finitely supported measures
through its coordinate projections. Secondly, we specify the convergence
requirements to be maintained throughout the inverse system and show
that they extend to each countable limit stage.

\subsection{Prediction of measure sequences}\label{subsec:prediction}

We shall apply diamond to sequences of finitely supported measures on $2^{\omega_1}$.  The union of the supports of one such sequence is countable.  The next observation ensures that sufficiently long coordinate projections preserve all its atoms, coefficients and variation norms.

\begin{lemma}\label{lem:support-projection}
Let $S\subseteq2^{\omega_1}$ be countable, and let
$p_\alpha:2^{\omega_1}\to2^\alpha$ be restriction of coordinates.
There is $\eta<\omega_1$ such that, whenever
$\eta\leqslant\alpha<\omega_1$, the map $p_\alpha$ is injective on $S$.
For every such $\alpha$, the map
\[
 (p_\alpha)_*:L(S)\longrightarrow L(p_\alpha[S]),
 \qquad \sum_{j<r}a_j\delta_{x_j}
       \longmapsto\sum_{j<r}a_j\delta_{p_\alpha(x_j)},
\]
is a linear isometric bijection, where $L(S)$ denotes the finitely
supported measures whose supports are contained in $S$; no closedness
of $S$ is required.
\end{lemma}
\begin{proof}
For every pair of distinct points $x,y\in S$, choose
$\xi(x,y)<\omega_1$ with $x(\xi(x,y))\neq y(\xi(x,y))$.
The set of these coordinates is countable.  Choose a countable
$\eta\geqslant\omega$ strictly greater than every chosen coordinate.
Then $p_\alpha(x)\neq p_\alpha(y)$ whenever $\alpha\geqslant\eta$.
Consequently, if $\mu=\sum_{j<r}a_j\delta_{x_j}$ is a reduced finite
representation with $x_j\in S$, then its push-forward has the same
nonzero coefficients at distinct points, and hence
\[
 \|(p_\alpha)_*\mu\|=\sum_{j<r}|a_j|=\|\mu\|.
\]
Every finitely supported measure on $p_\alpha[S]$ has a unique such
lift to $S$, proving the last assertion.
\end{proof}

We recall the set-theoretic terminology used below.  A subset of $\omega_1$ is \emph{club} if it is unbounded and is closed under suprema of increasing countable sequences whose supremum is below $\omega_1$.  A subset is \emph{stationary} if it meets every club.  The intersection of a stationary set with a club is stationary, because the intersection of two clubs is a club.  In particular, every stationary set meets every tail $[\eta,\omega_1)$.

Jensen's principle $\diamondsuit_{\omega_1}$ asserts the existence of a sequence $(d_\alpha)_{\alpha<\omega_1}$ with $d_\alpha\subseteq\alpha$ such that
\[
 \{\alpha<\omega_1:d_\alpha=E\cap\alpha\}
\]
is stationary for every $E\subseteq\omega_1$.  Thus diamond guesses the initial segments of each subset of $\omega_1$ at stationarily many stages.

The following lemma predicts the measure sequences directly. Related
forms of measure prediction appear in \cite[Lemma~6.4]{Glodkowski}
and \cite[Lemma~26]{KoszmiderSilber}. We give the elementary argument
for finitely supported measures.

\begin{lemma}\label{lem:diamond}
Assume $\diamondsuit_{\omega_1}$. There are sequences
$g^\alpha=(g_n^\alpha)_{n<\omega}$ in $L(2^\alpha)$, for
$\omega\leqslant\alpha<\omega_1$, such that, for every sequence
$(v_n)_{n<\omega}$ in $L(2^{\omega_1})$, stationarily many $\alpha$
satisfy
\[
 g_n^\alpha=(p_\alpha)_*v_n\qquad(n<\omega),
\]
and $p_\alpha$ is injective on $\bigcup_n\supp v_n$.
\end{lemma}
\begin{proof}
Fix bijections $e_2:\omega^2\to\omega$ and
$e_3:\omega^3\to\omega$, and enumerate $\mathbb Q$ as
$(q_k)_{k<\omega}$. The fixed points
\[
 C=\{\alpha<\omega_1:\alpha>\omega,\ \omega\cdot\alpha=\alpha\}
\]
form a club: unboundedness follows by countably iterating the map
$\xi\mapsto\omega\cdot\xi$ above any prescribed countable ordinal,
and closedness follows from its continuity at limits.

For $\delta\in C\cup\{\omega_1\}$, consider a labelled array
$\mathcal R=(r_n,a_{n,j},x_{n,j})$, where $r_n<\omega$,
$a_{n,j}\in\mathbb R$ and $x_{n,j}\in2^\delta$ for $j<r_n$.
Empty lists and repeated points are allowed. Represent this array by
\[
\begin{split}
 E_\delta(\mathcal R)={}&
 \{e_3(n,r_n,0):n<\omega\}\\
 &{}\cup\{e_3(n,j,k+1):n,k<\omega,\ j<r_n,\ q_k<a_{n,j}\}\\
 &{}\cup\{\omega\cdot(\xi+1)+e_2(n,j):n<\omega,\ j<r_n,
          \ \xi<\delta,\ x_{n,j}(\xi)=1\}.
\end{split}
\]
Here the last line uses ordinal arithmetic. The first two lines lie
in $\omega$ and recover all list lengths and coefficients, using
the rational lower cuts. Every entry in the last line lies outside
$\omega$, and ordinal division by $\omega$ recovers $\xi+1$ and
the label $(n,j)$. Thus the array is uniquely determined by its
representing subset of $\delta$.

Let $(d_\alpha)_{\alpha<\omega_1}$ witness diamond. If $\alpha\in C$
and $d_\alpha=E_\alpha(\mathcal R)$ for such an array, put
$g_n^\alpha=\sum_{j<r_n}a_{n,j}\delta_{x_{n,j}}$; otherwise let
$g_n^\alpha=0$ for every $n$. These are measures, so repeated atoms
are combined and their coefficients may cancel.

Given $(v_n)_n$, choose finite representations
$v_n=\sum_{j<r_n}a_{n,j}\delta_{x_{n,j}}$ and let $\mathcal R$ be
the resulting array on $2^{\omega_1}$. For $\alpha\in C$, an entry
$\omega\cdot(\xi+1)+m$ lies below $\alpha$ exactly when
$\xi<\alpha$. All entries describing lengths and coefficients also
lie below $\alpha$. Consequently
\[
 E_{\omega_1}(\mathcal R)\cap\alpha
   =E_\alpha(\mathcal R\restriction\alpha),
\]
where restriction leaves lengths and coefficients unchanged and
replaces each point by its first $\alpha$ coordinates. Diamond
therefore gives the asserted equality of measures at stationarily
many stages in $C$. Finally, intersect these stages with the tail
given by Lemma~\ref{lem:support-projection} for
$\bigcup_n\supp v_n$.
\end{proof}

For a stage space $K_\alpha\subseteq2^\alpha$, the prediction
$g^\alpha$ is \emph{valid} if every $g_n^\alpha$ belongs to
$L(K_\alpha)$ and
$A_\alpha=\Span\{g_n^\alpha:n<\omega\}$ is infinite dimensional.
Otherwise it is \emph{invalid}. In particular, a zero prediction
is invalid.

\subsection{Preservation at countable limits}\label{subsec:countable-limits}

We specify the convergence requirements for the inverse system to be constructed in Section~\ref{sec:recursion} and prove that they are preserved at countable limits.  For subsets $I,J\subseteq\N$, the notation $J\subseteq^*I$ means that $J\setminus I$ is finite.  Convergence along an infinite $J\subseteq\N$ always means convergence as the index tends to infinity through $J$.

At a valid prediction stage $\theta$, let $(D_{\theta,k})_{k<\omega}$
enumerate the clopen subsets of $K_\theta$. The construction introduces
one sequence of measures for each $D_{\theta,k}$.
We index these sequences by $i=(\theta,k)$ and write $b_i=\theta+1$
for their initial stage. All values of $k$ at the same $\theta$ refer
to one predicted subspace; the corresponding localised limits will be
used together in Section~\ref{sec:coupling}.
The sequence with index $i$ is
\[
 (a_{i,n})_{n<\omega}\subseteq L(K_{b_i})
\]
of norm-one measures with pairwise disjoint finite supports.  At every proper stage $\alpha\geqslant b_i$, we retain an infinite set $J_i^\alpha\subseteq\N$ and a measure $\mu_i^\alpha\in M(K_\alpha)$ with $\|\mu_i^\alpha\|\leqslant1$.  Define
\begin{equation}\label{eq:global-lift-sets}
 B_{i,n}^\alpha=
 \{\zeta\in M(K_\alpha):\|\zeta\|\leqslant1,
 (\pi_{b_i,\alpha})_*\zeta=a_{i,n}\}.
\end{equation}
These sets are nonempty and weak-star compact by Lemma~\ref{lem:lifts}.  The three preservation requirements are
\begin{align}
 J_i^\beta&\subseteq^*J_i^\alpha
 &&(b_i\leqslant\alpha\leqslant\beta<\omega_1),
 \label{eq:J-invariant}\\
 (\pi_{\alpha,\beta})_*\mu_i^\beta&=\mu_i^\alpha
 &&(b_i\leqslant\alpha\leqslant\beta<\omega_1),
 \label{eq:mu-invariant}\\
 \sup_{\zeta\in B_{i,n}^\alpha}
 |\zeta(f)-\mu_i^\alpha(f)|&\longrightarrow0
 &&(n\to\infty\text{ along }J_i^\alpha),\quad f\in C(K_\alpha).
 \label{eq:uniform-invariant}
\end{align}
At the initial stage $b_i$, the map $\pi_{b_i,b_i}$ is the identity, so $B_{i,n}^{b_i}=\{a_{i,n}\}$.  Consequently, if $a_{i,n}\xrightarrow{w^*}\mu_i^{b_i}$, the initial requirements hold with $J_i^{b_i}=\N$.

The next lemma allows the construction to continue at every countable limit ordinal.  It produces one infinite set $J_i^\rho$ almost contained in all the earlier sets $J_i^\alpha$, and a measure $\mu_i^\rho$ satisfying the same projection and convergence identities.

\begin{lemma}\label{lem:limit-invariant}
Let $\omega<\rho<\omega_1$ be a limit ordinal, and suppose that the inverse system and requirements \eqref{eq:J-invariant}--\eqref{eq:uniform-invariant} have been constructed below $\rho$.  With $K_\rho$ the inverse limit, all three requirements extend to $\rho$ for every index $i$ with $b_i<\rho$.
\end{lemma}
\begin{proof}
Fix an index $i$ with $b_i<\rho$, and choose an increasing sequence $(\alpha_m)_{m<\omega}$ cofinal in $\rho$, with $\alpha_0\geqslant b_i$.  Each finite intersection $\bigcap_{j\leqslant m}J_i^{\alpha_j}$ is infinite: by \eqref{eq:J-invariant}, it contains all but finitely many members of the infinite set $J_i^{\alpha_m}$.  Inductively choose
\[
 n_m\in\bigcap_{j\leqslant m}J_i^{\alpha_j},
 \qquad n_m>n_{m-1}\quad(m>0),
 \qquad J_i^\rho=\{n_m:m<\omega\}.
\]
For any $\alpha\in[b_i,\rho)$, choose $j$ with $\alpha\leqslant\alpha_j$.  All $n_m$ with $m\geqslant j$ lie in $J_i^{\alpha_j}$, and $J_i^{\alpha_j}\setminus J_i^\alpha$ is finite.  Thus $J_i^\rho\subseteq^*J_i^\alpha$, verifying \eqref{eq:J-invariant}.

Consider the subalgebra
\[
 \mathcal E_i^\rho=
 \bigcup_{b_i\leqslant\alpha<\rho}
 \{g\circ\pi_{\alpha,\rho}:g\in C(K_\alpha)\}
 \subseteq C(K_\rho).
\]
The spaces in this union are increasing with $\alpha$, so the union is a unital algebra.  The projections separate points of the inverse limit; hence this algebra separates points and is uniformly dense by the Stone--Weierstrass theorem.  Set
\[
 \Phi_i(g\circ\pi_{\alpha,\rho})=\mu_i^\alpha(g).
\]
To verify well-definedness, suppose
$g\circ\pi_{\alpha,\rho}=h\circ\pi_{\beta,\rho}$, and choose $\delta\in[\max\{\alpha,\beta\},\rho)$.  The surjectivity of $\pi_{\delta,\rho}$ gives
$g\circ\pi_{\alpha,\delta}=h\circ\pi_{\beta,\delta}$.  Applying $\mu_i^\delta$ and using \eqref{eq:mu-invariant} yields $\mu_i^\alpha(g)=\mu_i^\beta(h)$.  The same passage to a common stage proves linearity.  Surjectivity also gives
\[
 |\Phi_i(g\circ\pi_{\alpha,\rho})|
 \leqslant\|g\|_\infty
 =\|g\circ\pi_{\alpha,\rho}\|_\infty.
\]
Thus $\Phi_i$ extends uniquely to a bounded linear functional of norm at most one on $C(K_\rho)$.  The Riesz representation theorem supplies $\mu_i^\rho\in M(K_\rho)$ with $\|\mu_i^\rho\|\leqslant1$.  For every $\alpha\in[b_i,\rho)$ and $g\in C(K_\alpha)$, its definition gives
\[
 ((\pi_{\alpha,\rho})_*\mu_i^\rho)(g)
 =\mu_i^\rho(g\circ\pi_{\alpha,\rho})
 =\mu_i^\alpha(g).
\]
Equality on continuous functions is equality of Radon measures.  This verifies \eqref{eq:mu-invariant} at $\rho$.

Finally, fix $f\in C(K_\rho)$ and $\eps>0$.  Density supplies $\alpha\in[b_i,\rho)$ and $g\in C(K_\alpha)$ with
$\|f-g\circ\pi_{\alpha,\rho}\|_\infty<\eps$.
For any $\zeta\in B_{i,n}^\rho$, the push-forward
$\eta=(\pi_{\alpha,\rho})_*\zeta$ has norm at most one and satisfies
$(\pi_{b_i,\alpha})_*\eta=a_{i,n}$; hence $\eta\in B_{i,n}^\alpha$.  Therefore
\[
 \sup_{\zeta\in B_{i,n}^\rho}|\zeta(f)-\mu_i^\rho(f)|
 \leqslant 2\eps+
 \sup_{\eta\in B_{i,n}^\alpha}|\eta(g)-\mu_i^\alpha(g)|.
\]
The last term tends to zero along $J_i^\rho$, because it tends to zero along $J_i^\alpha$ and $J_i^\rho\subseteq^*J_i^\alpha$.  Letting $\eps$ decrease to zero proves \eqref{eq:uniform-invariant}.  For each index $i$ with $b_i<\rho$, the construction of $J_i^\rho$ uses only the sets $J_i^\alpha$ for that same index, and the construction of $\mu_i^\rho$ uses only the measures $\mu_i^\alpha$.  None of \eqref{eq:J-invariant}--\eqref{eq:uniform-invariant} imposes a condition relating two different indices.  There are at most countably many such indices, since they are pairs $(\theta,k)$ with $\theta<\rho$ and $k<\omega$.  Applying the preceding argument separately to each of them proves the simultaneous assertion.
\end{proof}

\section{The successor lemma and the inverse system}\label{sec:successor}\label{sec:recursion}

Our aim is to construct, under $\diamondsuit_{\omega_1}$, a continuous
inverse system with $K_\omega=2^\omega$ and $K_\alpha\subseteq2^\alpha$
satisfying \eqref{eq:J-invariant}--\eqref{eq:uniform-invariant}.
The successor lemma below is proved in ZFC. It constructs one extension
that preserves the preceding uniform limits and introduces new
clopen-localised limits.

Diamond is used in Subsection~\ref{subsec:recursion} to obtain the
prediction sequences from Lemma~\ref{lem:diamond}.
At a valid prediction, Lemma~\ref{lem:flat} supplies the signed sequence
to which the successor lemma applies; its extension becomes
$K_{\alpha+1}$ and its limits give \eqref{eq:birth-limit}.
Lemma~\ref{lem:limit-invariant} handles countable limit stages.
The resulting weak-star closure statement is Lemma~\ref{lem:final-closure}.

\subsection{The successor extension}\label{subsec:successor}

\begin{lemma}\label{lem:simultaneous-successor}
Let $K$ be compact, metrisable and zero-dimensional, and let $I$ be a
finite or countably infinite initial segment of $\N$, possibly empty.
For $i\in I$, let $X_i$ be compact Hausdorff, let $r_i:K\to X_i$ be a
continuous surjection, and let $(a_{i,n})_{n<\omega}\subseteq L(X_i)$
have norm one and pairwise disjoint supports. Put
\[
 E_{i,n}=r_i^{-1}(\supp a_{i,n}),\qquad
 B_{i,n}=\{\zeta\in M(K):\|\zeta\|\leqslant1,
                 (r_i)_*\zeta=a_{i,n}\}.
\]
Let $J_i\subseteq\N$ be infinite and suppose that $\mu_i\in M(K)$ satisfies
\begin{equation}\label{eq:old-uniform-hypothesis}
 \sup_{\zeta\in B_{i,n}}|\zeta(f)-\mu_i(f)|\longrightarrow0
 \quad(n\to\infty,\ n\in J_i),\qquad f\in C(K).
\end{equation}
Let $(\sigma_n)_{n<\omega}\subseteq L(K)$ have pairwise disjoint supports
and Jordan decompositions $\sigma_n=p_n-q_n$, where
\[
 p_n,q_n\in P(K),\qquad
 p_n\xrightarrow{w^*}\lambda,\quad q_n\xrightarrow{w^*}\lambda
 \quad\text{for some }\lambda\in P(K).
\]
Thus $\|\sigma_n\|=2$ and $\sigma_n\to0$ weak-star.
Enumerate the clopen subsets of $K$ as $(D_k)_{k<\omega}$, including $K$,
and fix a surjection $\kappa:\N\to\N$ with all fibres infinite.

Then there exist a compact set
$G\subseteq K$ with $\lambda(G)>0$, a set
$F=\supp(\lambda\restr{G})\subseteq\supp\lambda$, and an open set
$U\subseteq K$ such that $\lambda(F)=\lambda(G)>0$,
$U\cap F=\varnothing$ and $\overline U=U\cup F$.
If $\lambda\restr{G}$ is nonatomic, then $F$ is perfect and uncountable;
if $\lambda\restr{G}=c\delta_x$ for some $x\in K$ and $c>0$, then
$F=\{x\}$.
The space
\[
 K'=((K\setminus U)\times\{0\})\cup((U\cup F)\times\{1\})
\]
is compact, metrisable and zero-dimensional. Its projection $\pi:K'\to K$
has the Borel section $s(x)=(x,\one_U(x))$. The following hold.
\begin{enumerate}
\item For each $i\in I$, there is an infinite $J'_i\subseteq J_i$ such that
\begin{equation}\label{eq:successor-old-conclusion}
 \sup_{\substack{\zeta'\in M(K'),\ \|\zeta'\|\leqslant1\\
                  (r_i\circ\pi)_*\zeta'=a_{i,n}}}
 |\zeta'(h)-(s_*\mu_i)(h)|\longrightarrow0
 \quad(n\to\infty,\ n\in J'_i),\qquad h\in C(K').
\end{equation}
Moreover, $\pi_*s_*\mu_i=\mu_i$ and $\|s_*\mu_i\|=\|\mu_i\|$.
\item There are strictly increasing indices $t_m$ such that
$\supp\sigma_{t_m}$ misses $F$. Its unique norm-preserving lift
$\widehat\sigma_{t_m}$ to $K'$ is finitely supported, and
\begin{equation}\label{eq:fresh-localised-limit}
 \frac12\widehat\sigma_{t_m}\xrightarrow{w^*}\nu_k
 \quad(m\to\infty,\ \kappa(m)=k),\qquad
 \nu_k=\frac12\int_{F\cap D_k}
       (\delta_{(x,1)}-\delta_{(x,0)})\,d\lambda(x).
\end{equation}
The measures $\widehat\sigma_{t_m}/2$ have norm one and pairwise disjoint
supports.
\end{enumerate}
\end{lemma}

\begin{figure}[H]
\centering
\begin{tikzpicture}[x=0.9cm,y=0.82cm,font=\small]
\draw[rounded corners] (0,0) rectangle (13,5.9);
\node[anchor=north west] at (0.12,5.8) {$K$};
\filldraw[fill=blue!3,draw=blue!65!black,dashed,rounded corners]
 (4.8,0.55) rectangle (12.5,5.2);
\node[anchor=north east] at (12.35,5.1) {$C_m$};
\filldraw[fill=gray!25,draw=black] (8.6,2.55) rectangle (11.7,3.3);
\node at (10.15,2.93) {$F$};
\draw[fill=blue!12,rounded corners] (5.25,3.45) rectangle (7.9,4.75);
\node at (6.57,4.4) {$U_m\supseteq T_m$};
\node at (5.9,3.86) {$+$};\node at (7.15,3.86) {$+$};
\draw[fill=orange!12,rounded corners] (5.25,0.95) rectangle (7.9,2.25);
\node at (6.57,1.92) {$L_m$};
\node at (5.9,1.35) {$-$};\node at (7.15,1.35) {$+$};
\draw[fill=green!8,rounded corners] (0.6,0.95) rectangle (3.7,2.4);
\node at (2.15,2.05) {$W_{i,m}$};
\node at (2.15,1.4) {$E_{i,n_{i,m}}$};
\draw[fill=blue!12,rounded corners] (0.6,3.45) rectangle (3.7,4.75);
\node at (2.15,4.4) {$P_m$};
\node at (1.5,3.86) {$+$};\node at (2.8,3.86) {$-$};
\node[align=center] at (10.15,1.38)
 {later $C_\ell$ avoid\\$U_m,L_m,W_{i,m}$};
\end{tikzpicture}
\caption{One step of the construction in the proof of Lemma~\ref{lem:simultaneous-successor}. The signs indicate points of $\supp p_{t_m}$ and $\supp q_{t_m}$. The points of $\supp\sigma_{t_m}$ in $P_m\cup T_m$ have second coordinate $1$ in $K'$; all remaining points of this support lie in $L_m$ and have second coordinate $0$. The set $W_{i,m}$ contains $E_{i,n_{i,m}}$ and is disjoint from every later $C_\ell$. Every point of $F$ has two preimages. The figure is schematic: $L_m$ may extend outside $C_m$, and $P_m$ may meet $W_{i,m}$.}
\label{fig:successor}
\end{figure}

\begin{proof}
By Lemma~\ref{lem:lifts}, every $B_{i,n}$ is nonempty and weak-star
compact, and its members have norm one and are concentrated on
$E_{i,n}$. For fixed $i$, the sets $E_{i,n}$ are pairwise disjoint and compact. Choosing any member of $B_{i,n}$ for $n\in J_i$ and using
\eqref{eq:old-uniform-hypothesis} gives
$|\mu_i(f)|\leqslant\|f\|$ for every $f\in C(K)$; hence
$\|\mu_i\|\leqslant1$.

Add the family $E_{-1,n}=\supp\sigma_n$. For each
$i\in I\cup\{-1\}$, disjointness of the sets $E_{i,n}$ and finiteness
of $\lambda$ imply that their tail unions have measures tending to zero.
Choose $N_i$ with
\[
 \lambda\left(\bigcup_{n\geqslant N_i}E_{i,n}\right)<2^{-i-4}.
\]
Since $\sum_{i=-1}^{\infty}2^{-i-4}=1/4$, the union of all these tail
unions has measure less than $1/4$. By inner regularity, its complement
contains a compact set $G$ of positive measure. Take
$F=\supp(\lambda\restr{G})$. The measure $\lambda\restr{G}$ is
concentrated on its support, so $F\subseteq G$ is compact and
$\lambda\restr{F}=\lambda\restr{G}$. In particular,
\begin{equation}\label{eq:F-properties}
 \lambda(F)>0,\qquad F=\supp(\lambda\restr{F}),\qquad
 F\cap E_{i,n}=\varnothing
 \quad(i\in I\cup\{-1\},\ n\geqslant N_i).
\end{equation}
The inclusion $F\subseteq\supp\lambda$ follows because
$\lambda\restr{G}\leqslant\lambda$. Suppose that $\lambda\restr{G}$
is nonatomic. If $x$ were isolated in $F$, some open set $V$ would satisfy
$V\cap F=\{x\}$. The support condition and concentration on $F$ would
then give
\[
 0<(\lambda\restr{G})(V)=(\lambda\restr{G})(\{x\}),
\]
contrary to nonatomicity. Thus $F$ is perfect. If it were countable,
its closed singletons would express it as a countable union of nowhere
dense sets, contradicting the Baire category theorem. Hence $F$ is
uncountable. If $\lambda\restr{G}=c\delta_x$ for $c>0$, its support
is $\{x\}$. The construction also applies when the restriction has
several atoms or both atomic and nonatomic parts.

Fix a compatible metric $d$ on $K$ and a norm-dense sequence
$(f_j)_{j<\omega}$ in $C(K)$ containing every clopen indicator.
Such a sequence exists because a compact metrisable zero-dimensional
space has only countably many clopen sets: each is a finite union of
members of a fixed countable clopen base.
We recursively choose clopen sets $U_m,L_m,C_m$, clopen neighbourhoods
$W_{i,m}$ for $i\in I$, $i\leqslant m$, and indices $t_m,n_{i,m}$.
Write $P_m=\bigcup_{\ell<m}U_\ell$; thus $P_0=\varnothing$.
All the sets $U_m,L_m,W_{i,m}$ will miss $F$, whereas $C_m$ will be a
clopen neighbourhood of $F$.

At step $m$, for each $i\in I$ with $i\leqslant m$, choose
$n_{i,m}\in J_i$, with $n_{i,m}\geqslant N_i$, larger than the indices
previously chosen for this $i$, and satisfying
\begin{equation}\label{eq:old-rank-estimates}
 \sup_{\zeta\in B_{i,n_{i,m}}}|\zeta(g)-\mu_i(g)|<2^{-m}
 \quad\bigl(g\in\{f_j,f_j\one_{P_m}:j\leqslant m\}\bigr).
\end{equation}
Indeed, $P_m$ has already been chosen and is clopen, so all the
functions $f_j$ and $f_j\one_{P_m}$ with $j\leqslant m$ are fixed
members of $C(K)$. For each such function $g$,
\eqref{eq:old-uniform-hypothesis} gives an integer $N(g)$ such that
\[
 \sup_{\zeta\in B_{i,n}}|\zeta(g)-\mu_i(g)|<2^{-m}
 \qquad(n\in J_i,\ n\geqslant N(g)).
\]
There are finitely many $g$. An element of the infinite set $J_i$
larger than all their thresholds, at least $N_i$, and larger than the
previously selected indices therefore satisfies
\eqref{eq:old-rank-estimates}.
The compact set $E_{i,n_{i,m}}$ misses $F$ by
\eqref{eq:F-properties}. Zero-dimensionality and compactness give a
clopen neighbourhood $W_{i,m}$ of $E_{i,n_{i,m}}$ disjoint from $F$.

Let $Q_m$ be the union of all previously chosen $U_\ell,L_\ell$
($\ell<m$) and all $W_{i,\ell}$ chosen up to and including step $m$.
This is a finite union of clopen sets disjoint from $F$.
Since $\lambda$ is Radon,
$\lambda(F)=\inf\{\lambda(O):F\subseteq O,\ O\text{ open in }K\}$.
We may therefore choose an open neighbourhood $O_m$ of $F$ with
$\lambda(O_m\setminus F)<2^{-m}$. Choose a clopen neighbourhood $C_m$
of $F$ contained in
\[
 O_m\cap(K\setminus Q_m)\cap\{x:d(x,F)<2^{-m}\},
\]
and also contained in $C_{m-1}$ if $m>0$. For $m=0$ there is no
additional condition involving a previous $C_m$. Each displayed open
set contains the compact set $F$, so the clopen neighbourhood exists
by a finite subcover from a clopen base. We have
\begin{equation}\label{eq:C-rank-estimates}
 C_m\subseteq\{x:d(x,F)<2^{-m}\},\qquad
 \lambda(C_m\setminus F)<2^{-m}.
\end{equation}

Set $D=D_{\kappa(m)}$. Choose $t_m\geqslant N_{-1}$, larger than all
earlier $t_\ell$, so that, for $j\leqslant m$,
\begin{align}
 |\sigma_{t_m}(f_j)|&<2^{-m},\label{eq:fresh-rank1}\\
 |\sigma_{t_m}(f_j\one_{P_m})|&<2^{-m},\label{eq:fresh-rank2}\\
 |p_{t_m}(f_j\one_{C_m\cap D})
       -\lambda(f_j\one_{C_m\cap D})|&<2^{-m}.
       \label{eq:fresh-rank3}
\end{align}
To justify this common choice of $t_m$, all the sets $P_m$, $C_m$ and
$D$ have been fixed and are clopen. For each $j\leqslant m$, weak-star
convergence gives
\[
 \sigma_n(f_j)\longrightarrow0,\qquad
 \sigma_n(f_j\one_{P_m})\longrightarrow0,\qquad
 p_n(f_j\one_{C_m\cap D})
   \longrightarrow\lambda(f_j\one_{C_m\cap D}).
\]
Each of these finitely many convergences gives a threshold beyond which
the corresponding error is less than $2^{-m}$. Taking $t_m$ larger
than all these thresholds, all earlier $t_\ell$, and at least $N_{-1}$
gives \eqref{eq:fresh-rank1}--\eqref{eq:fresh-rank3}.

Put $T_m=\supp p_{t_m}\cap C_m\cap D$. This finite set misses $F$.
Choose a clopen $U_m\subseteq C_m\cap D$ containing $T_m$, disjoint from
$F$, and meeting $\supp\sigma_{t_m}$ in exactly $T_m$.
For each point of $T_m$, a clopen neighbourhood inside $C_m\cap D$
can be chosen to avoid the compact set
$F\cup(\supp\sigma_{t_m}\setminus T_m)$; the union of these finitely
many neighbourhoods works. Take $U_m=\varnothing$ when $T_m$ is empty.
Since $C_m$ misses every earlier $U_\ell$, the sets $U_m$ are pairwise
disjoint. The finite set
\[
 \supp\sigma_{t_m}\setminus(P_m\cup T_m)
\]
is disjoint from the compact set $F\cup P_m\cup U_m$. Choose a clopen
neighbourhood $L_m$ of that finite set avoiding this compact set.
This finishes step $m$. The construction ensures that all later
$C_\ell$ and $U_\ell$ miss $L_m$ and every $W_{i,m}$.


Let $U=\bigcup_mU_m$, an open set disjoint from $F$. We prove
$\overline U\setminus U=F$. If $y\in\overline U\setminus U$, choose
$y_j\in U$ with $y_j\to y$ and write $y_j\in U_{m_j}$.
Each finite union of the clopen sets $U_m$ is closed and excludes $y$;
thus $m_j\to\infty$. Equation~\eqref{eq:C-rank-estimates} gives
$d(y_j,F)<2^{-m_j}\to0$, so $y\in F$.
Conversely, let $x\in F$ and let $V$ be a neighbourhood of $x$.
Choose a clopen neighbourhood $V_0$ with $x\in V_0\subseteq V$.
The support condition in \eqref{eq:F-properties} gives
$\lambda(F\cap V_0)>0$. The function $\one_{V_0}$ is one of the $f_j$.
Infinitely many steps have $D_{\kappa(m)}=K$, and at all sufficiently
large such steps \eqref{eq:fresh-rank3} yields
\[
 p_{t_m}(C_m\cap V_0)
 >\lambda(C_m\cap V_0)-2^{-m}
 \geqslant\lambda(F\cap V_0)-2^{-m}>0.
\]
Consequently $T_m\cap V_0\neq\varnothing$, and $U_m$ meets $V$.
This proves $F\subseteq\overline U$ and the asserted identity.

Put
\[
 K'_0=(K\setminus U)\times\{0\},\qquad
 K'_1=(U\cup F)\times\{1\}.
\]
Both sets are closed in $K\times\{0,1\}$, because $K\setminus U$
and $U\cup F=\overline U$ are closed. Thus $K'=K'_0\cup K'_1$ is
compact, metrisable and zero-dimensional. The sets $K'_0$ and $K'_1$
are also clopen in $K'$, since
$K'_e=K'\cap(K\times\{e\})$ for $e=0,1$. The function
\[
 \chi_1:K'\longrightarrow\R,\qquad
 \chi_1(x,e)=e=\one_{K'_1}(x,e),
\]
is therefore continuous. The projection $\pi(x,e)=x$ is a continuous
surjection, with
\[
 \pi^{-1}\{x\}=
 \begin{cases}
 \{(x,1)\},&x\in U,\\
 \{(x,0),(x,1)\},&x\in F,\\
 \{(x,0)\},&x\in K\setminus(U\cup F).
 \end{cases}
\]
As $U$ is Borel, $s(x)=(x,\one_U(x))$ is a Borel map into $K'$
and $\pi\circ s=\mathrm{id}_K$.
For a signed Radon measure $\mu$ on $K$, define
$(s_*\mu)(B)=\mu(s^{-1}(B))$ for every Borel set $B\subseteq K'$.
This finite Borel measure is Radon because $K'$ is compact and
metrisable. For $h\in C(K')$ it satisfies
$(s_*\mu)(h)=\int_K h(s(x))\,d\mu(x)$, and
\[
 \pi_*s_*\mu=\mu,\qquad
 \|\mu\|=\|\pi_*s_*\mu\|\leqslant\|s_*\mu\|
          \leqslant\|\mu\|.
\]
The first equality follows from $\pi\circ s=\mathrm{id}_K$.
For either map $T=\pi$ or $T=s$ and every continuous function $h$
on its target, the integral formula gives
\[
 |(T_*\eta)(h)|=\left|\int h\circ T\,d\eta\right|
 \leqslant\|h\|\,\|\eta\|.
\]
Taking the supremum over $\|h\|\leqslant1$ gives the two norm
inequalities.

We determine the intersections of $U$ with the supports of the measures.
Fix $i\in I$ and $m\geqslant i$. Since $E_{i,n_{i,m}}\subseteq W_{i,m}$
and $C_\ell\cap W_{i,m}=\varnothing$ for every $\ell\geqslant m$,
none of the sets $U_\ell\subseteq C_\ell$ with $\ell\geqslant m$
meets $E_{i,n_{i,m}}$. As $P_m=\bigcup_{\ell<m}U_\ell$, we obtain
\begin{equation}\label{eq:old-selector-identity}
 \one_U\restr{E_{i,n_{i,m}}}
 =\one_{P_m}\restr{E_{i,n_{i,m}}}.
\end{equation}
Next put $S_m=\supp\sigma_{t_m}$. By construction,
$U_m\cap S_m=T_m$ and $S_m\setminus(P_m\cup T_m)\subseteq L_m$.
Since $T_m\subseteq U_m$, this gives
$S_m\subseteq P_m\cup U_m\cup L_m$.
For every $\ell>m$, the choice of $C_\ell$ ensures that
$U_\ell\subseteq C_\ell$ is disjoint from $P_m\cup U_m\cup L_m$,
and hence from $S_m$. Thus only the sets $U_\ell$ with $\ell\leqslant m$
contribute to $U\cap S_m$, and
\[
 U\cap S_m=(P_m\cap S_m)\cup T_m.
\]
The two sets on the right are disjoint, because
$T_m\subseteq C_m$ and $C_m\cap P_m=\varnothing$.
The support $S_m$ misses $F$, so each $x\in S_m$ has the unique
preimage $s(x)$. Writing $\sigma_{t_m}=\sum_{x\in S_m}c_{m,x}\delta_x$,
we obtain the norm-preserving lift
\[
 \widehat\sigma_{t_m}=\sum_{x\in S_m}c_{m,x}\delta_{s(x)},
 \qquad \|\widehat\sigma_{t_m}\|=
 \sum_{x\in S_m}|c_{m,x}|=2.
\]
By Lemma~\ref{lem:lifts}, applied after division by two, every
norm-preserving lift is concentrated on $\pi^{-1}(S_m)$.
Since each of these fibres is a singleton, the push-forward requirement
fixes its coefficient at $s(x)$ to be $c_{m,x}$. This proves uniqueness.
Since $T_m\subseteq\supp p_{t_m}$, the disjointness of the Jordan
supports gives, for every $g\in C(K)$,
\begin{equation}\label{eq:fresh-selector-identity}
 \widehat\sigma_{t_m}((g\circ\pi)\chi_1)
 =\sigma_{t_m}(g\one_{P_m})
       +p_{t_m}(g\one_{C_m\cap D_{\kappa(m)}}).
\end{equation}
Indeed, $\sigma_{t_m}\restr{T_m}=p_{t_m}\restr{T_m}$ and
$p_{t_m}(g\one_{T_m})=p_{t_m}(g\one_{C_m\cap D_{\kappa(m)}})$.
The measures $\widehat\sigma_{t_m}$ have pairwise disjoint supports,
because their supports project onto the pairwise disjoint sets $S_m$.

Fix $i\in I$ and put $J'_i=\{n_{i,m}:m\geqslant i\}$.
The map $m\mapsto n_{i,m}$ is strictly increasing, so convergence
as $m\to\infty$ is the same as convergence as $n\to\infty$ along $J'_i$.
Given $h\in C(K')$, the function $x\mapsto h(x,0)$ is continuous
on the closed set $K\setminus U$, and $x\mapsto h(x,1)$ is continuous
on the closed set $U\cup F$. By the Tietze extension theorem, choose
$h_0,h_1\in C(K)$ such that
\[
 h_0(x)=h(x,0)\quad(x\in K\setminus U),\qquad
 h_1(x)=h(x,1)\quad(x\in U\cup F).
\]
Let $\zeta'$ be any measure occurring in the supremum in
\eqref{eq:successor-old-conclusion} with $n=n_{i,m}$, and set
$\zeta=\pi_*\zeta'$. Then $\zeta\in B_{i,n_{i,m}}$.
Write $E=E_{i,n_{i,m}}$. Lemma~\ref{lem:lifts} for
$r_i\circ\pi$ gives
\[
 |\zeta'|(K'\setminus\pi^{-1}(E))=0,
 \qquad |\zeta|(K\setminus E)=0.
\]
Since $E\cap F=\varnothing$, the map
$\pi:\pi^{-1}(E)\to E$ is a continuous bijection between compact
Hausdorff spaces, hence a homeomorphism, whose inverse is $s\restr{E}$.
For a Borel set $B\subseteq K'$, this implies
\[
 \begin{split}
 \zeta'(B)
 &=\zeta'(B\cap\pi^{-1}(E))\\
 &=\zeta'\bigl(\pi^{-1}(s^{-1}(B))\cap\pi^{-1}(E)\bigr)
 =\zeta(s^{-1}(B)).
 \end{split}
\]
Thus $\zeta'=s_*\zeta$. Moreover, \eqref{eq:old-selector-identity}
says that $s(x)=(x,\one_{P_m}(x))$ for $x\in E$, so
\[
 \zeta'(h)=\zeta\bigl(h_0+(h_1-h_0)\one_{P_m}\bigr).
\]
Choose fixed indices $j_0,j_1$ with $\|h_e-f_{j_e}\|<\eps$ for $e=0,1$.
The function $h_0+(h_1-h_0)\one_{P_m}$ differs by less than $\eps$
from $f_{j_0}+(f_{j_1}-f_{j_0})\one_{P_m}$: the difference is
$h_0-f_{j_0}$ off $P_m$ and $h_1-f_{j_1}$ on $P_m$.
For $m\geqslant\max(i,j_0,j_1)$, the three estimates in
\eqref{eq:old-rank-estimates} and the norm bounds thus give
\[
 \left|\zeta'(h)-
 \mu_i\bigl(h_0+(h_1-h_0)\one_{P_m}\bigr)\right|
 \leqslant2\eps+3\cdot2^{-m},
\]
uniformly over all these $\zeta'$.
Since $P_m$ increases to $U$, dominated convergence with respect to
$|\mu_i|$ gives
\[
 \mu_i\bigl(h_0+(h_1-h_0)\one_{P_m}\bigr)
 \longrightarrow\mu_i\bigl(h_0+(h_1-h_0)\one_U\bigr)
 =(s_*\mu_i)(h).
\]
Letting first $m\to\infty$ and then $\eps\to0$ proves part (1).

For part (2), fix $k$ and restrict to $m\in\kappa^{-1}\{k\}$.
For each fixed $j$ and all $m\geqslant j$ in this set,
\eqref{eq:fresh-rank1} gives
$|\widehat\sigma_{t_m}(f_j\circ\pi)|<2^{-m}$.
Equations \eqref{eq:C-rank-estimates},
\eqref{eq:fresh-rank2}--\eqref{eq:fresh-rank3} and
\eqref{eq:fresh-selector-identity} give the further estimate
\[
 \left|\widehat\sigma_{t_m}((f_j\circ\pi)\chi_1)
       -\int_{F\cap D_k}f_j\,d\lambda\right|
 \leqslant(2+\|f_j\|)2^{-m}.
\]
For $g\in C(K)$ with $\|g-f_j\|<\eps$, the left side with $g$
in place of $f_j$ is at most
$3\eps+(2+\|f_j\|)2^{-m}$: the lifted measure has norm two,
and $g\mapsto\int_{F\cap D_k}g\,d\lambda$ is a continuous linear
functional of norm at most one. Similarly,
$|\widehat\sigma_{t_m}(g\circ\pi)|\leqslant2\eps+2^{-m}$.
These estimates extend both limits to every $g\in C(K)$.
For an arbitrary $h\in C(K')$, choose the extensions $h_0,h_1$
as in the proof of part (1). Then
\[
 h=h_0\circ\pi+((h_1-h_0)\circ\pi)\chi_1.
\]
The two established limits yield
\[
 \widehat\sigma_{t_m}(h)\longrightarrow
 \int_{F\cap D_k}\bigl(h(x,1)-h(x,0)\bigr)\,d\lambda(x).
\]
Division by two proves \eqref{eq:fresh-localised-limit} and completes
the proof.
\end{proof}

\subsection{The inverse system and its final measures}\label{subsec:recursion}

Assume $\diamondsuit_{\omega_1}$ and fix the sequences $g^\alpha$
supplied by Lemma~\ref{lem:diamond}. Fix also one map
$\kappa:\N\to\N$ with every fibre infinite, and enumerate each fibre as
\[
 \kappa^{-1}\{k\}=\{m_k(0)<m_k(1)<\cdots\}\qquad(k<\omega).
\]
We use this same map at every successor stage.
We construct a continuous inverse system
\[
 (K_\alpha,\pi_{\alpha,\beta})_{
 \omega\leqslant\alpha\leqslant\beta\leqslant\omega_1},
 \qquad K_\alpha\subseteq2^\alpha,
\]
whose bonding maps are coordinate projections.  Start with $K_\omega=2^\omega$.  Write $\mathcal I_\alpha$ for the set of indices $i$ of measure sequences introduced at or before stage $\alpha$.  Initially $\mathcal I_\omega=\varnothing$: no sequences $(a_{i,n})_n$, index sets $J_i^\omega$ or measures $\mu_i^\omega$ have yet been chosen, so the requirements \eqref{eq:J-invariant}--\eqref{eq:uniform-invariant} are vacuous.  At each limit stage take the inverse limit.  Each successor map described below is onto.  At a limit, compactness and the finite intersection property show that every point of every earlier space extends to a point of the inverse limit.  Thus all bonding maps are surjective.

Inductively every $K_\alpha$ is compact and zero-dimensional: this holds at the start, is preserved by the successor lemma and the constant-coordinate extension, and is preserved by inverse limits. Every proper stage is metrisable, being a subspace of $2^\alpha$ for countable $\alpha$.

Suppose the system and all the measure data have been constructed through a proper stage $\alpha$.  Each earlier successor introduced at most countably many sequences, and $\alpha$ is countable, so $\mathcal I_\alpha$ is countable.  We distinguish the two cases according to the validity of $g^\alpha$, as defined in Subsection~\ref{subsec:prediction}.

If the prediction is invalid, put $K_{\alpha+1}=K_\alpha\times\{0\}\subseteq2^{\alpha+1}$, with constant new coordinate, and let $s_\alpha(x)=(x,0)$.  For each $i\in\mathcal I_\alpha$, set
\[
 J_i^{\alpha+1}=J_i^\alpha,
 \qquad \mu_i^{\alpha+1}=(s_\alpha)_*\mu_i^\alpha.
\]
The map $\pi_{\alpha,\alpha+1}$ is a homeomorphism with inverse $s_\alpha$, so its push-forward identifies the corresponding lift sets and all three preservation requirements hold.  No new sequences are introduced, and $\mathcal I_{\alpha+1}=\mathcal I_\alpha$.

For a valid prediction, apply Lemma~\ref{lem:flat} to $A_\alpha$, obtaining disjointly supported measures
\[
 \sigma_n=p_n-q_n\in A_\alpha,
 \qquad p_n,q_n\xrightarrow{w^*}\lambda_\alpha,
 \qquad\|\sigma_n\|=2.
\]
Enumerate the clopen subsets of $K_\alpha$ without repetitions as $(D_{\alpha,k})_{k<\omega}$.  This is possible because $K_\alpha$ is compact metrisable and zero-dimensional; its clopen algebra is countable and infinite, since $K_\alpha$ maps onto $2^\omega$.  Apply Lemma~\ref{lem:simultaneous-successor} with this enumeration and the fixed map $\kappa$.  For each $i\in\mathcal I_\alpha$, the required data are
\[
 r_i=\pi_{b_i,\alpha},\quad
 (a_{i,n})_{n<\omega},\quad J_i^\alpha,\quad\mu_i^\alpha.
\]
The compact sets used by that lemma are
\[
 E_{i,n}^\alpha=
 \pi_{b_i,\alpha}^{-1}(\supp a_{i,n});
\]
for each fixed $i$ they are pairwise disjoint.  If $\mathcal I_\alpha$ is finite or empty, the lemma is applied to precisely this finite or empty family of data.

For this application enumerate $\mathcal I_\alpha$ by an initial segment of $\N$, as in the hypotheses of Lemma~\ref{lem:simultaneous-successor}, and then use the original indices $i\in\mathcal I_\alpha$ for the resulting data.

Let $F_\alpha,U_\alpha$ be the sets furnished by the lemma.  Identify the space $K'\subseteq K_\alpha\times\{0,1\}$ from that lemma with $K_{\alpha+1}\subseteq2^{\alpha+1}$, using its second coordinate as coordinate $\alpha$, and put
\[
 s_\alpha(x)=(x,\one_{U_\alpha}(x)).
\]
For each $i\in\mathcal I_\alpha$, take $J_i^{\alpha+1}$ to be the infinite subset of $J_i^\alpha$ supplied by the lemma, and set $\mu_i^{\alpha+1}=(s_\alpha)_*\mu_i^\alpha$.  We check the requirements separately.  First,
$J_i^{\alpha+1}\subseteq J_i^\alpha$, which together with the earlier almost inclusions proves \eqref{eq:J-invariant}.  Secondly,
$\pi_{\alpha,\alpha+1}\circ s_\alpha=\mathrm{id}_{K_\alpha}$, so
$(\pi_{\alpha,\alpha+1})_*\mu_i^{\alpha+1}=\mu_i^\alpha$; the earlier coherence identities then give \eqref{eq:mu-invariant} at every earlier stage.  Thirdly, part~(1) of Lemma~\ref{lem:simultaneous-successor} is exactly \eqref{eq:uniform-invariant} at $\alpha+1$.  The norm bound is retained because Borel push-forward does not increase total variation.

The successor lemma supplies distinct indices $t_m$ and their unique coefficient-preserving lifts $\widehat\sigma_{t_m}\in L(K_{\alpha+1})$.  For the new index $i=(\alpha,k)$, put $b_i=\alpha+1$ and introduce the sequence
\begin{equation}\label{eq:birth-sequence}
 a_{(\alpha,k),n}=
 \frac12\widehat\sigma_{t_{m_k(n)}},
 \qquad n<\omega.
\end{equation}
This is a sequence indexed by $\N$, of norm-one measures with pairwise disjoint finite supports.  Distinct values of $k$ use disjoint sets of indices $m$, so no selected measure occurs in two of the newly introduced sequences.  Set $J_i^{b_i}=\N$ and
\begin{equation}\label{eq:birth-limit}
 \mu_{(\alpha,k)}^{\alpha+1}
 =\frac12\int_{F_\alpha\cap D_{\alpha,k}}
  (\delta_{(x,1)}-\delta_{(x,0)})\,d\lambda_\alpha(x).
\end{equation}
Part~(2) of the successor lemma gives convergence of \eqref{eq:birth-sequence} to \eqref{eq:birth-limit}.  Since $B_{i,n}^{b_i}=\{a_{i,n}\}$, this verifies \eqref{eq:uniform-invariant}.  The other two requirements are identities at the initial stage $b_i$.  A zero initial limit is allowed and requires no special case.  Set $\mathcal I_{\alpha+1}=\mathcal I_\alpha\cup(\{\alpha\}\times\N)$.

At a countable limit $\rho$, put $\mathcal I_\rho=\bigcup_{\omega\leqslant\alpha<\rho}\mathcal I_\alpha$ and apply Lemma~\ref{lem:limit-invariant} to each $i\in\mathcal I_\rho$.  This completes the construction through all proper stages.  Put $K=K_{\omega_1}$; it is a compact zero-dimensional space mapping onto $2^\omega$, and hence infinite, with weight at most $\aleph_1$.

For each index $i$ introduced in the construction, the coherent measures determine a final measure $\mu_i\in M(K)$ directly from the exact factorisation in Lemma~\ref{lem:factor}.  If
$f=g\circ\pi_{\alpha,\omega_1}$ with $\alpha\geqslant b_i$, define
\begin{equation}\label{eq:final-record-measure}
 \mu_i(f)=\mu_i^\alpha(g).
\end{equation}
If two such representations are given, pass to a common proper stage and use surjectivity and \eqref{eq:mu-invariant}, exactly as in the well-definedness argument of Lemma~\ref{lem:limit-invariant}.  The values coincide.  Passing to a common stage also proves linearity, and
$|\mu_i(f)|\leqslant\|f\|_\infty$ follows from surjectivity.  Thus \eqref{eq:final-record-measure} defines a bounded functional on all of $C(K)$, and the Riesz representation theorem identifies it with a signed Radon measure of norm at most one.  Its projection to every $K_\alpha$, $\alpha\geqslant b_i$, is $\mu_i^\alpha$.  There is no choice of a set $J_i^{\omega_1}$ in the construction.

The uniform convergence requirement was imposed for all norm-preserving lifts so that it applies to any measures on the final space having the prescribed initial projections.  The next lemma gives the resulting weak-star closure statement.

\begin{lemma}\label{lem:final-closure}
Let $i$ be an index introduced in the construction and let $A\subseteq L(K)$ be a linear subspace.  If $A$ contains norm-one measures $(w_n)_{n<\omega}$ with
\[
 (\pi_{b_i,\omega_1})_*w_n=a_{i,n}\qquad(n<\omega),
\]
then $\mu_i\in\overline A^{\,w^*}$, where the closure is taken in $M(K)$.
\end{lemma}
\begin{proof}
Let $f_1,\ldots,f_r\in C(K)$ and $\eps>0$ specify a basic weak-star neighbourhood of $\mu_i$.  Lemma~\ref{lem:factor} gives one proper stage $\beta\geqslant b_i$ and functions $g_j\in C(K_\beta)$ with
$f_j=g_j\circ\pi_{\beta,\omega_1}$ for $1\leqslant j\leqslant r$.  The measure
$(\pi_{\beta,\omega_1})_*w_n$ belongs to $B_{i,n}^\beta$ for every $n$.  By \eqref{eq:uniform-invariant}, all sufficiently large $n$ in the infinite set $J_i^\beta$ satisfy, simultaneously for $1\leqslant j\leqslant r$,
\[
 |w_n(f_j)-\mu_i(f_j)|
 =|((\pi_{\beta,\omega_1})_*w_n)(g_j)-\mu_i^\beta(g_j)|
 <\eps.
\]
Thus every such neighbourhood meets $A$.
\end{proof}

\section{From localised limits to differences of point masses}\label{sec:coupling}

For one valid prediction, we represent all its clopen-localised limits
using a single positive Radon measure $\gamma$ on $K^2$, supported on
distinct pairs with the same projection to the predicted stage.
We prove that, if every localised limit belongs to the weak-star closure
of a subspace $A\subseteq L(K)$, then every
$(y_0,y_1)\in\supp\gamma$ satisfies
$\delta_{y_1}-\delta_{y_0}\in\overline A^{\,w^*}$.
Since $\gamma\neq0$, this gives the nonzero finitely supported measure
needed to contradict relative closedness of $A$ in the next section.

Fix a valid prediction at a stage $\theta$, and put $b=\theta+1$, $F=F_\theta$, $\lambda=\lambda_\theta$, and $\lambda_F=\lambda\restr{F}$.  The measure $\lambda_F$ is finite, positive and nonzero.  For a clopen $D\subseteq K_\theta$, let $k$ be its unique index in $(D_{\theta,k})_{k<\omega}$ and use the explicit abbreviations
\[
 \mu_D^\beta=\mu_{(\theta,k)}^\beta
 \quad(b\leqslant\beta<\omega_1),
 \qquad \mu_D=\mu_{(\theta,k)}.
\]
Thus $\mu_D^\beta$ and $\mu_D$ are the limits associated with the sequence indexed by $(\theta,k)$.

For $e\in\{0,1\}$ and $b\leqslant\beta<\omega_1$, we construct Borel maps
\[
 t_e^\beta:F\longrightarrow K_\beta
\]
with coherent projections.  At stage $b$ set $t_e^b(x)=(x,e)$; this is continuous.  At a successor set
$t_e^{\beta+1}=s_\beta\circ t_e^\beta$, which is Borel because both maps are Borel.  Since
$\pi_{\beta,\beta+1}\circ s_\beta=\mathrm{id}$, this definition retains coherence with $t_e^\beta$ and hence with all earlier maps.  At a countable limit $\rho$, define $t_e^\rho(x)$ as the unique point of the inverse limit with the prescribed coherent earlier projections.  Its coordinate at each $\xi<\rho$ is the corresponding coordinate of some $t_e^\beta$ with $b\leqslant\beta<\rho$ and $\xi<\beta$, and is therefore Borel.  Since $\rho$ is countable, $2^\rho$ has a countable base of finite-coordinate cylinders.  Their inverse images are Borel, so $t_e^\rho$ is Borel as a map into $K_\rho$.  Coherence holds by the inverse-limit definition.  The induction also gives
\begin{equation}\label{eq:section-coherence}
 \pi_{\beta,\delta}\circ t_e^\delta=t_e^\beta
 \quad(b\leqslant\beta\leqslant\delta<\omega_1),
 \qquad \pi_{\theta,\beta}(t_e^\beta(x))=x.
\end{equation}

At every proper stage the localisation formula is the following identity of Radon measures:
\begin{equation}\label{eq:localisation-proper}
 \mu_D^\beta=
 \frac12\bigl((t_1^\beta)_*-(t_0^\beta)_*\bigr)
 (\lambda\restr{F\cap D}).
\end{equation}
All push-forwards on the right are finite Borel measures on the compact metrisable space $K_\beta$ and are therefore Radon.  At stage $b$ the identity is \eqref{eq:birth-limit}.  If it holds at $\beta$, push both sides forward by the same Borel map $s_\beta$.  The composition rule for Borel push-forwards and the definition of $t_e^{\beta+1}$ give the identity at $\beta+1$.  This is legitimate for bounded Borel functions: equality of Radon measures entails equality of their integrals against such functions, even though $f\circ s_\beta$ need not be continuous.

At a countable limit $\rho$, let the right-hand side of \eqref{eq:localisation-proper} define $\nu_D^\rho$.  For every $\beta\in[b,\rho)$, coherence \eqref{eq:section-coherence} gives
$(\pi_{\beta,\rho})_*\nu_D^\rho=\mu_D^\beta$.
The same holds for $\mu_D^\rho$ by \eqref{eq:mu-invariant}.  Hence these two measures agree on every continuous function pulled back from a stage below $\rho$.  These functions are uniformly dense by the Stone--Weierstrass argument in Lemma~\ref{lem:limit-invariant}; both measures are bounded, so they agree on all of $C(K_\rho)$.  This proves \eqref{eq:localisation-proper} at every proper stage.

Define the positive Radon measures
\[
 \gamma^\beta=(t_0^\beta,t_1^\beta)_*\lambda_F
 \quad\text{on }K_\beta^2,
 \qquad b\leqslant\beta<\omega_1.
\]
Their mass is $\lambda(F)$, and \eqref{eq:section-coherence} makes them coherent under the pair projections
$Q_{\beta,\delta}=\pi_{\beta,\delta}\times\pi_{\beta,\delta}$.
To obtain their final extension, let $Q_\beta=\pi_{\beta,\omega_1}\times\pi_{\beta,\omega_1}$ and consider
\[
 \mathcal E=
 \bigcup_{b\leqslant\beta<\omega_1}
 \{h\circ Q_\beta:h\in C(K_\beta^2)\}
 \subseteq C(K^2).
\]
This is a unital algebra.  The pair projections separate points of $K^2$, so the Stone--Weierstrass theorem makes $\mathcal E$ uniformly dense.  Set
\[
 \Phi(h\circ Q_\beta)=\gamma^\beta(h).
\]
If two representations define the same function, pass to a common proper stage.  Surjectivity of the pair projection then gives equality of the two functions at that stage, and coherence of the $\gamma^\beta$ gives equality of their integrals.  Thus $\Phi$ is well defined and linear.  Surjectivity also gives
\[
 |\Phi(h\circ Q_\beta)|
 \leqslant\lambda(F)\|h\circ Q_\beta\|_\infty,
 \qquad \Phi(\one)=\lambda(F).
\]
If $h\circ Q_\beta\geqslant0$, then $h\geqslant0$, so $\Phi$ is positive.  It therefore extends to a positive bounded functional on $C(K^2)$; positivity follows, for example, by approximating a nonnegative function uniformly by members of $\mathcal E$ and adding their approximation errors times $\one$.  Since $K^2$ is compact Hausdorff, the Riesz representation theorem yields a positive Radon measure $\gamma$ with
\begin{equation}\label{eq:coupling-marginals}
 (Q_\beta)_*\gamma=\gamma^\beta,
 \qquad \gamma(K^2)=\lambda(F)>0.
\end{equation}

The set
\[
 R^b=\{((x,0),(x,1)):x\in F\}\subseteq K_b^2
\]
is compact and carries $\gamma^b$.  Its closed preimage is
\begin{equation}\label{eq:initial-relation}
 \begin{split}
 R=Q_b^{-1}(R^b)
 =\{(y_0,y_1)\in K^2:
 &\pi_{\theta,\omega_1}(y_0)=\pi_{\theta,\omega_1}(y_1)\in F,\\
 &(y_0)_\theta=0,\ (y_1)_\theta=1\}.
 \end{split}
\end{equation}
By \eqref{eq:coupling-marginals},
$\gamma(K^2\setminus R)=\gamma^b(K_b^2\setminus R^b)=0$.
As $R$ is closed, $\supp\gamma\subseteq R$, and $R$ is disjoint from the diagonal.  Write $p_0,p_1:K^2\to K$ for the two coordinate projections and let $r:R\to F$ be their common projection to $K_\theta$.

For a clopen $D\subseteq K_\theta$, the weight
\[
 \chi_D=\one_D\circ\pi_{\theta,\omega_1}\circ p_0
\]
is continuous on $K^2$ and agrees with $\one_{r^{-1}(D)}$ on $R$.  The final localisation identity is
\begin{equation}\label{eq:localisation-final}
 \mu_D=
 \frac12\bigl((p_1)_*-(p_0)_*\bigr)(\chi_D\gamma).
\end{equation}
Indeed, given $f\in C(K)$, choose $\beta\geqslant b$ and $\bar f\in C(K_\beta)$ with $f=\bar f\circ\pi_{\beta,\omega_1}$.  The continuous function
\[
 (y_0,y_1)\longmapsto
 \chi_D(y_0,y_1)\bigl(f(y_1)-f(y_0)\bigr)
\]
factors through $Q_\beta$.  Its integral against $\gamma$ therefore equals the integral of its factor against $\gamma^\beta$.  By the definition of $\gamma^\beta$, one half of this integral is
\[
 \frac12\int_{F\cap D}
  \bigl(\bar f(t_1^\beta(x))-\bar f(t_0^\beta(x))\bigr)
 \,d\lambda(x)
 =\mu_D^\beta(\bar f)=\mu_D(f).
\]
This proves \eqref{eq:localisation-final} on continuous functions and hence as an identity of Radon measures.

The following lemma converts the closure condition for the localised measures into a pointwise equality on $\supp\gamma$.

\begin{lemma}\label{lem:coupling}
Let $A\subseteq L(K)$ be a linear subspace such that
$\mu_D\in\overline A^{\,w^*}$ for every clopen $D\subseteq K_\theta$, with closure in $M(K)$.  Then
\begin{equation}\label{eq:common-support-equalities}
 \supp\gamma\subseteq
 \bigcap_{f\in A_\perp}
 \{(y_0,y_1)\in K^2:f(y_0)=f(y_1)\}.
\end{equation}
Consequently every $(y_0,y_1)\in\supp\gamma$ satisfies
$\delta_{y_1}-\delta_{y_0}\in\overline A^{\,w^*}\cap L(K)$.
\end{lemma}
\begin{proof}
Fix $f\in A_\perp$ and write $f=\bar f\circ\pi_{\beta,\omega_1}$ for a proper $\beta\geqslant b$.  Every measure in $\overline A^{\,w^*}$ annihilates $f$, since evaluation at $f$ is weak-star continuous.  Equations \eqref{eq:localisation-proper} and \eqref{eq:final-record-measure} therefore yield
\[
 \int_{F\cap D}h\,d\lambda=0
 \quad(D\subseteq K_\theta\text{ clopen}),
 \qquad h=\bar f\circ t_1^\beta-\bar f\circ t_0^\beta.
\]
The function $h$ is bounded and Borel on $F$.  The traces on $F$ of the clopen subsets of $K_\theta$ form an algebra generating the Borel sets of $F$, because $K_\theta$ is zero-dimensional and has a countable clopen base.  The finite signed measure $h\lambda_F$ vanishes on that algebra, including $F$ itself; uniqueness of finite measures on a generating algebra, or the monotone class theorem, gives $h\lambda_F=0$.  It follows that $h=0$ almost everywhere for $\lambda_F$.

By \eqref{eq:coupling-marginals} and the definition of $\gamma^\beta$,
\[
 \int_{K^2}|f(y_1)-f(y_0)|^2\,d\gamma(y_0,y_1)
 =\int_F|h|^2\,d\lambda=0.
\]
If this continuous nonnegative integrand were positive at a point of $\supp\gamma$, it would be bounded below by a positive constant on some neighbourhood of that point; the neighbourhood has positive $\gamma$-measure, contradicting the zero integral.  It therefore vanishes throughout $\supp\gamma$.  This proves containment in the closed equality set for the fixed $f$.  Since the same fixed support is contained in that set for every $f\in A_\perp$, \eqref{eq:common-support-equalities} follows.

For each pair in this support, its difference of point masses annihilates $A_\perp$.  The annihilator identity \eqref{eq:bipolar}, applied in $(M(K),C(K))$, puts it in the closure of $A$ in $M(K)$.  It is finitely supported, so it also lies in $L(K)$, as asserted.
\end{proof}

\section{Proof of Theorem A}\label{sec:finish}

We first prove the quotient assertions of Theorem~A.

\begin{proof}[Proof of the quotient assertions]
Let $K$ be the compactum constructed above.  By Theorem~B and Lemma~\ref{lem:dual}, the quotient assertions follow once every countably infinite-dimensional subspace of $L(K)$ is shown not to be relatively weak-star closed.  Each such subspace will be treated at one correctly predicted stage of the construction.

Fix
\[
 A=\Span\{v_n:n<\omega\}\subseteq L(K),
 \qquad\dim A=\aleph_0,
 \qquad S=\bigcup_{n<\omega}\supp v_n.
\]
Here $L(S)$ denotes the measures in $L(K)$ whose finite supports are contained in $S$; the countable set $S$ need not be closed.  In particular, $A\subseteq L(S)$.  Let us regard the $v_n$ as measures on $2^{\omega_1}$. Lemma~\ref{lem:diamond} gives a correct prediction stage $\theta$ at which $\pi_{\theta,\omega_1}$ is injective on $S$.

Distinct atoms of any measure in $L(S)$ remain distinct under this projection.  Thus push-forward is injective and variation-isometric on $L(S)$.  In particular, the predicted span at stage $\theta$ is
\[
 A_\theta=
 \Span\{(\pi_{\theta,\omega_1})_*v_n:n<\omega\}
 =(\pi_{\theta,\omega_1})_*A,
\]
and it is infinite dimensional.  Its measures are supported in $K_\theta$, so this is a valid prediction.  The restriction of push-forward to $A$ is consequently a linear isometric bijection from $A$ onto $A_\theta$.

Let $(\sigma_n)$ be the sequence chosen from $A_\theta$ at this successor step.  There are unique $w_n\in A$ with
\[
 (\pi_{\theta,\omega_1})_*w_n=\sigma_n,
 \qquad \|w_n\|=\|\sigma_n\|=2.
\]
Their supports are pairwise disjoint: an atom belonging to two of them would project to an atom in the two corresponding disjoint supports of the $\sigma_n$.  Put $b=\theta+1$.  For each selected index $t_m$, define
\[
 \tau_m=(\pi_{b,\omega_1})_*w_{t_m}.
\]
Composition of the projections and their norm contractivity give
\begin{equation}\label{eq:norm-preserving-hinge}
 2=\|\sigma_{t_m}\|
 =\|(\pi_{\theta,b})_*\tau_m\|
 \leqslant\|\tau_m\|
 \leqslant\|w_{t_m}\|=2.
\end{equation}
Hence $\tau_m/2$ is a norm-preserving lift of $\sigma_{t_m}/2$.  Lemma~\ref{lem:lifts} shows that it is concentrated on
$\pi_{\theta,b}^{-1}(\supp\sigma_{t_m})$.
The successor construction makes every fibre over $\supp\sigma_{t_m}$ a singleton, because that support misses $F_\theta$.  Thus the projection on this finite preimage is a bijection, and its prescribed push-forward fixes the coefficient at every point.  It follows that
\begin{equation}\label{eq:identified-lift}
 (\pi_{b,\omega_1})_*w_{t_m}
 =\tau_m=\widehat\sigma_{t_m}.
\end{equation}

Fix any $k<\omega$.  Reindexing the norm-one final measures
\[
 \left(\frac12w_{t_{m_k(n)}}\right)_{n<\omega}
 \subseteq A
\]
according to \eqref{eq:birth-sequence}, identity \eqref{eq:identified-lift} shows that their projections to $K_b$ are exactly the sequence $(a_{(\theta,k),n})_{n<\omega}$.  Lemma~\ref{lem:final-closure} now gives
\[
 \mu_D\in\overline A^{\,w^*}
 \qquad(D\subseteq K_\theta\text{ clopen}).
\]

Take the measure $\gamma$ constructed in Section~\ref{sec:coupling} for this prediction.  Its mass is $\lambda_\theta(F_\theta)>0$, so its support is nonempty.  Choose $(y_0,y_1)\in\supp\gamma$.  Lemma~\ref{lem:coupling} gives
\[
 \eta=\delta_{y_1}-\delta_{y_0}
 \in\overline A^{\,w^*}\cap L(K).
\]
Because $\supp\gamma\subseteq R$, relation \eqref{eq:initial-relation} implies $y_0\neq y_1$ and
$\pi_{\theta,\omega_1}(y_0)=\pi_{\theta,\omega_1}(y_1)$.
Consequently $\eta\neq0$ whereas $(\pi_{\theta,\omega_1})_*\eta=0$.  Push-forward is injective on $A\subseteq L(S)$, so $\eta\notin A$.  This proves that $A$ is not relatively weak-star closed.

Lemma~\ref{lem:dual} excludes infinite-dimensional metrisable quotients, and Theorem~B excludes infinite-dimensional separable quotients.
\end{proof}

The remaining assertions of Theorem~A follow from the topological properties proved in the next section.

\section{Further properties of the compact space}\label{sec:properties}

We first verify that the successor maps used above preserve enough of the topology of the initial Cantor space to give separability and crowdedness at the final stage.

\begin{lemma}\label{lem:irreducible-successor}
Every projection produced by Lemma~\ref{lem:simultaneous-successor} is irreducible: no proper closed subset of its domain maps onto its range.
\end{lemma}
\begin{proof}
Retain the notation of that lemma. We already know $F\subseteq\overline U$. We claim also that
\begin{equation}\label{eq:lower-boundary}
F\subseteq\overline{K\setminus(U\cup F)}.
\end{equation}
Fix $x\in F$ and a clopen neighbourhood $V$ of $x$. Then $\lambda(F\cap V)>0$ by \eqref{eq:F-properties}. Consider sufficiently large steps $m$ for which $D_{\kappa(m)}=K$ and $\one_V$ occurs among $f_0,\ldots,f_m$. Since $C_m\cap P_m=\varnothing$, equations \eqref{eq:fresh-rank1}--\eqref{eq:fresh-rank3} give
\begin{align*}
q_{t_m}(V\setminus P_m)
&=p_{t_m}(V\setminus P_m)-\sigma_{t_m}(V)
                                  +\sigma_{t_m}(V\cap P_m)\\
&\geqslant p_{t_m}(C_m\cap V)-2\cdot2^{-m}\\
&\geqslant\lambda(F\cap V)-3\cdot2^{-m}>0.
\end{align*}
Choose a point of $\supp q_{t_m}$ in $V\setminus P_m$. It does not belong to $T_m\subseteq\supp p_{t_m}$, so it belongs to $L_m$. By construction, $L_m\cap F=\varnothing$ and $L_m\cap U_j=\varnothing$ for every $j<\omega$. Our point therefore lies in $V\setminus(U\cup F)$, proving \eqref{eq:lower-boundary}.

The singleton fibres of $\pi$ are those over $K\setminus F$. They are dense in $K'$: the points of $U\times\{1\}$ are dense at every $(x,1)$ with $x\in F$ because $F\subseteq\overline U$, and the points of $(K\setminus(U\cup F))\times\{0\}$ are dense at every $(x,0)$ with $x\in F$ by \eqref{eq:lower-boundary}. A closed subset of $K'$ mapping onto $K$ must contain every singleton fibre and hence their closure $K'$. This proves irreducibility.
\end{proof}

Irreducibility passes to the limit and transfers separability and crowdedness from the initial Cantor space.

\begin{proposition}\label{prop:separable-crowded}
The space $K$ constructed above is separable and crowded.
\end{proposition}
\begin{proof}
For a continuous surjection $r:X\to Y$ between compact Hausdorff spaces, irreducibility is equivalent to the following condition: every nonempty open $V\subseteq X$ contains $r^{-1}(W)$ for some nonempty open $W\subseteq Y$. Indeed, take $W=Y\setminus r[X\setminus V]$; it is open by compactness and nonempty precisely when the proper closed set $X\setminus V$ fails to map onto $Y$. This condition also shows that a composition of irreducible maps is irreducible.

We prove by transfinite induction that $\pi_{\omega,\rho}:K_\rho\to2^\omega$ is irreducible for every $\omega\leqslant\rho\leqslant\omega_1$. The assertion is immediate at $\rho=\omega$. At a successor, the new bonding map is irreducible by Lemma~\ref{lem:irreducible-successor}, or is a homeomorphism in the invalid case, and composition applies. At a limit $\rho$, a nonempty open subset of $K_\rho$ contains a nonempty cylinder $\pi_{\beta,\rho}^{-1}(V)$ for some $\omega\leqslant\beta<\rho$ and nonempty open $V\subseteq K_\beta$. The inductive hypothesis supplies a nonempty open $W\subseteq2^\omega$ with $\pi_{\omega,\beta}^{-1}(W)\subseteq V$. Consequently $\pi_{\omega,\rho}^{-1}(W)$ lies in the original open set, as required. This argument includes the final limit $\rho=\omega_1$.

Put $r=\pi_{\omega,\omega_1}$. Choose a countable dense subset $D\subseteq2^\omega$ and one point $x_d\in r^{-1}(d)$ for each $d\in D$. The open-set criterion just proved shows that every nonempty open subset of $K$ contains some $x_d$, so these points form a countable dense set. If $x\in K$ were isolated, that criterion applied to $\{x\}$ would give a nonempty open $W\subseteq2^\omega$ with $r^{-1}(W)\subseteq\{x\}$. Surjectivity would force $W=\{r(x)\}$, contradicting the absence of isolated points in the Cantor space.
\end{proof}

The quotient conclusion also excludes the two subspaces appearing in the definition of an Efimov space.

\begin{corollary}\label{cor:efimov}
The space $K$ is an Efimov space and has weight exactly $\aleph_1$.
\end{corollary}
\begin{proof}
An infinite compact space containing a nontrivial convergent sequence or a copy of $\beta\N$ has an infinite-dimensional metrisable quotient of its pointwise function space, by \cite[Corollary~3]{BKS2018}. The conclusion proved in Section~\ref{sec:finish} therefore excludes both subspaces of $K$. An infinite compact metrisable space contains a nontrivial convergent sequence, so $K$ is not metrisable. Since $K\subseteq2^{\omega_1}$, its weight is at most $\aleph_1$; a compact Hausdorff space of countable weight is metrisable. Thus its weight is $\aleph_1$.
\end{proof}

The absence of convergent sequences determines the cardinality through a binary tree of closed $G_\delta$ sets.

\begin{proposition}\label{prop:cardinality}
The cardinality of $K$ is $2^{\aleph_1}$.
\end{proposition}
\begin{proof}
There are no $G_\delta$ points in $K$. To see this, suppose $\{x\}=\bigcap_n O_n$ with each $O_n$ open. By regularity choose open neighbourhoods $V_n$ of $x$ such that $\overline V_0\subseteq O_0$ and $\overline V_{n+1}\subseteq V_n\cap O_{n+1}$. Compactness shows that $(V_n)$ is a local base at $x$: if $O$ is a neighbourhood of $x$ and no $\overline V_n$ is contained in $O$, the nested nonempty compact sets $\overline V_n\setminus O$ have an intersection, contradicting $\bigcap_n\overline V_n=\{x\}$. In a crowded Hausdorff space every nonempty open set is infinite. We can therefore choose pairwise distinct $x_n\in V_n\setminus\{x\}$, giving a nontrivial convergent sequence, contrary to Corollary~\ref{cor:efimov}.

It follows that every nonempty closed $G_\delta$ subset of $K$ has at least two points. Construct nonempty closed $G_\delta$ sets $F_s$, for $s\in2^{<\omega_1}$, starting with $F_\varnothing=K$. This is a set-sized recursion: $|2^{<\omega_1}|=2^{\aleph_0}$, since each countable level has cardinality at most $2^{\aleph_0}$, there are $\aleph_1\leqslant2^{\aleph_0}$ levels, and the level $2^\omega$ has cardinality $2^{\aleph_0}$. Given $F_s$, choose distinct $x_0,x_1\in F_s$ and a continuous $u:K\to[0,1]$ with $u(x_e)=e$. Set
\[
F_{s^\frown0}=F_s\cap u^{-1}([0,1/3]),\qquad
F_{s^\frown1}=F_s\cap u^{-1}([2/3,1]).
\]
These are disjoint nonempty closed $G_\delta$ sets: inverse images of the displayed closed intervals are zero sets and hence $G_\delta$. At a nonzero countable limit $\rho$, put $F_s=\bigcap_{\xi<\rho}F_{s\restriction\xi}$ for $s\in2^\rho$. This is a nonempty closed $G_\delta$ set by compactness and countability, so the recursion continues. For each branch $z\in2^{\omega_1}$, compactness gives a point in $\bigcap_{\xi<\omega_1}F_{z\restriction\xi}$. Distinct branches give disjoint intersections, yielding $2^{\aleph_1}$ distinct points of $K$. The reverse inequality follows from $K\subseteq2^{\omega_1}$.
\end{proof}

The classical theorem of Josefson and Nissenzweig \cite{Josefson,Nissenzweig} supplies norm-one weak-star null sequences in the dual of every infinite-dimensional Banach space; finite support is an additional requirement. Write
\[
\ell_1(X)=\left\{\sum_{x\in S}a_x\delta_x:
 S\subseteq X\text{ countable},\ \sum_{x\in S}|a_x|<\infty\right\}
\subseteq M(X).
\]
The space $C(X)$ has the \emph{$\ell_1$-Grothendieck property} if every weak-star convergent sequence in $\ell_1(X)$ is weakly convergent in $M(X)$.

\begin{corollary}\label{cor:fsjn}
The space $K$ admits no finitely supported Josefson--Nissenzweig sequence, and $C(K)$ has the $\ell_1$-Grothendieck property.
\end{corollary}
\begin{proof}
By \cite[Theorem~1]{BKS2019}, the existence of such a sequence is equivalent to $C_p(K)$ having a quotient isomorphic to $(c_0)_p$, where $(c_0)_p$ is $c_0$ with the topology inherited from $\R^\N$. This would be an infinite-dimensional Hausdorff metrisable quotient. The second assertion follows from \cite[Theorem~4.1]{KSZ}.
\end{proof}

The full Grothendieck property of $C(K)$ and the possible existence of an isomorphic copy of $\ell_\infty$ in $C(K)$ are not decided by this argument.

\section{An obstruction to one-point resolutions}\label{sec:obstruction}

The following elementary example explains why a single double fibre need not disturb the convergence of a sequence of differences of probability measures.

\begin{proposition}\label{prop:one-point}
On the Cantor space $X=2^\omega$ there are finitely supported probabilities $p_n,q_n$ such that all their supports are pairwise disjoint and $p_n,q_n\to\lambda$ weak-star for the nonatomic Bernoulli probability $\lambda$. If $\pi:Y\to X$ is a continuous surjection from a compact metrisable space and all its fibres are singletons except possibly one double fibre, then every choice of lifts $\widehat p_n,\widehat q_n$ obtained by placing each atom over its original point and retaining its coefficient satisfies
\[
\frac12(\widehat p_n-\widehat q_n)\xrightarrow{w^*}0,
\qquad
\left\|\frac12(\widehat p_n-\widehat q_n)\right\|=1.
\]
In particular, $Y$ admits a finitely supported Josefson--Nissenzweig sequence.
\end{proposition}
\begin{proof}
For each $n$ and each binary word $s$ of length $n$, choose two points $x^0_{n,s},x^1_{n,s}$ in its basic cylinder $[s]$, choosing every point distinct from all points chosen earlier. This is possible since every cylinder is infinite and at each step only finitely many points have been used. Set
\[
p_n=2^{-n}\sum_{s\in2^n}\delta_{x^0_{n,s}},\qquad
q_n=2^{-n}\sum_{s\in2^n}\delta_{x^1_{n,s}}.
\]
For $f\in C(X)$, both $|p_n(f)-\lambda(f)|$ and $|q_n(f)-\lambda(f)|$ are bounded by $\max_{s\in2^n}\operatorname{osc}(f,[s])$, which tends to zero by uniform continuity. This proves the asserted convergence.

Let $x$ be the exceptional base point, choosing any $x\in X$ if there is no exceptional fibre. We show that $\lambda$ has a unique probability lift to $Y$. On $X\setminus\{x\}$ the inverse of the singleton-fibre map is continuous. Indeed, if $y\neq x$ and an open set $O\subseteq Y$ contains the unique point over $y$, then $X\setminus\pi[Y\setminus O]$ is an open neighbourhood of $y$ whose full preimage lies in $O$. Thus for every $h\in C(Y)$, the function $g(y)=h(\pi^{-1}(y))$ on $X\setminus\{x\}$ extends, by assigning any value at $x$, to a bounded Borel function on $X$. Any probability $\nu$ with $\pi_*\nu=\lambda$ gives mass zero to $\pi^{-1}(x)$, since $\lambda(\{x\})=0$, and therefore
\[
\nu(h)=\int_X g\,d\lambda.
\]
This identity determines $\nu$ uniquely on $C(Y)$. Existence follows, for example, by choosing finite probability lifts of the $p_n$ and taking a weak-star cluster point in the compact space $P(Y)$.

For any such choices, every cluster point of $(\widehat p_n)$ and of $(\widehat q_n)$ is a probability projecting to $\lambda$. Both sequences therefore have the same unique cluster point. Since $P(Y)$ is compact metrisable, each sequence converges to that point: otherwise a subsequence outside a fixed neighbourhood would have a different cluster point. Their difference tends to zero. Disjointness of the original supports implies disjointness of all lifted supports, so each difference has norm two, as required.
\end{proof}

The same probability measures show directly that stationarily many successor maps in our construction double a perfect uncountable set.

\begin{proposition}\label{prop:perfect-splitting-stages}
In the inverse system of Section~\ref{sec:recursion}, stationarily many valid stages $\theta<\omega_1$ have nonatomic $\lambda_\theta$ and perfect, uncountable $F_\theta$.
\end{proposition}
\begin{proof}
Let $\rho$ be the Bernoulli probability on $K_\omega=2^\omega$ and choose finitely supported probabilities $p_j^0,q_j^0\to\rho$ with all supports pairwise disjoint, as in Proposition~\ref{prop:one-point}. Lift each atom once to $K=K_{\omega_1}$ and retain its coefficient, obtaining probabilities $p_j,q_j\in L(K)$. Put $v_j=p_j-q_j$ and $A=\Span\{v_j:j<\omega\}$. The initial projection is injective on the union of all these supports. Consequently, for every $\omega\leqslant\theta<\omega_1$, the measures
\[
p_j^\theta=(\pi_{\theta,\omega_1})_*p_j,\qquad
q_j^\theta=(\pi_{\theta,\omega_1})_*q_j
\]
have pairwise disjoint supports, and $v_j^\theta=p_j^\theta-q_j^\theta$ has norm two.

Lemma~\ref{lem:diamond} supplies stationarily many stages at which the predicted subspace is
$A_\theta=\Span\{v_j^\theta:j<\omega\}$. The vectors $v_j^\theta$ are linearly independent by disjointness of their supports, so every such prediction is valid. Fix one of these stages and write the sequence selected by Lemma~\ref{lem:flat} as
\[
\sigma_n=\sum_jc_{n,j}v_j^\theta,\qquad
I_n=\{j:c_{n,j}\neq0\}.
\]
The sums are finite, and disjointness of all block supports gives
$2=\|\sigma_n\|=2\sum_j|c_{n,j}|$.
Furthermore, the finite sets $I_n$ are pairwise disjoint. Indeed, if $c_{n,j}\neq0$, the entire support of $v_j^\theta$ is contained in $\supp\sigma_n$, because no other block can cancel any of its coefficients. The supports of the $\sigma_n$ are pairwise disjoint. Thus $I_n$ eventually avoids every finite subset of $\N$.

For $c\in\R$, set $c^+=\max\{c,0\}$ and $c^-=\max\{-c,0\}$. The Jordan parts satisfy
\[
\begin{split}
(\pi_{\omega,\theta})_*\sigma_n^+
 &=\sum_j(c_{n,j}^+p_j^0+c_{n,j}^-q_j^0),\\
(\pi_{\omega,\theta})_*\sigma_n^-
 &=\sum_j(c_{n,j}^+q_j^0+c_{n,j}^-p_j^0).
\end{split}
\]
These are convex combinations with total coefficient one. For $f\in C(2^\omega)$ and $N<\omega$, put
\[
\eps_N(f)=\sup_{j\geqslant N}
\max\{|p_j^0(f)-\rho(f)|,\ |q_j^0(f)-\rho(f)|\}.
\]
We have $\eps_N(f)\to0$. When $I_n\subseteq[N,\infty)$, the absolute difference between the evaluation of either displayed probability at $f$ and $\rho(f)$ is at most
$\sum_j|c_{n,j}|\eps_N(f)=\eps_N(f)$.
Hence both projected Jordan sequences converge weak-star to $\rho$. Their limits at stage $\theta$ are $\lambda_\theta$, so
$(\pi_{\omega,\theta})_*\lambda_\theta=\rho$.

The measure $\lambda_\theta$ is nonatomic. Indeed, the inverse images of the length-$r$ Cantor cylinders partition $K_\theta$ into finitely many sets of measure $2^{-r}$. For any Borel set of positive measure, choose $r$ with $2^{-r}$ smaller than that measure. One member of this partition cuts off a strictly smaller positive measure, ruling out an atom. The nonzero restriction of $\lambda_\theta$ used in Lemma~\ref{lem:simultaneous-successor} is therefore nonatomic, and that lemma gives a perfect, uncountable $F_\theta$. This holds at every correct prediction stage chosen above.
\end{proof}

The theorem on simple extensions recalled in Section~\ref{sec:prelim}, together with Corollary~\ref{cor:fsjn}, precludes a presentation of $K$ by a continuous inverse system starting with a singleton and using simple extensions at every successor. Proposition~\ref{prop:perfect-splitting-stages} identifies the difference directly in our system: stationarily many successor maps have uncountably many two-point fibres. At each such stage the map is also not fully closed, since the images of the disjoint closed sets $K_{\theta+1}\cap(K_\theta\times\{0\})$ and $K_{\theta+1}\cap(K_\theta\times\{1\})$ intersect in the infinite set $F_\theta$, whereas full closedness requires this intersection to be finite.

Fedorchuk's classical diamond example \cite{Fedorchuk} has the same topological properties listed in Theorem~A, but has the finitely supported Josefson--Nissenzweig property \cite[Corollary~6.12]{KSZ}. Our example therefore differs in its behaviour for finitely supported measures. Talagrand's compactum \cite{Talagrand}, whose $C(K)$ has the Grothendieck property, also admits no presentation by simple extensions, by \cite[Theorems~4.1 and 6.10]{KSZ}.

\section*{Acknowledgements}
The first-named author acknowledges the institutional support of the Mathematical Institute of the Czech Academy of Sciences (RVO: 67985840).

\begin{thebibliography}{99}

\bibitem{Arhangelskii}
A.~V.~Arkhangel'skii,
\emph{Topological Function Spaces},
Mathematics and its Applications \textbf{78}, Kluwer, Dordrecht, 1992.
\url{https://doi.org/10.1007/978-94-011-2598-7}

\bibitem{BKS2018}
T.~Banakh, J.~K\k{a}kol and W.~\'Sliwa,
\emph{Metrizable quotients of $C_p$-spaces},
Topology Appl. \textbf{249} (2018), 95--102.
\url{https://doi.org/10.1016/j.topol.2018.09.012}

\bibitem{BKS2019}
T.~Banakh, J.~K\k{a}kol and W.~\'Sliwa,
\emph{Josefson--Nissenzweig property for $C_p$-spaces},
Rev. R. Acad. Cienc. Exactas F\'{i}s. Nat. Ser. A Mat. RACSAM
\textbf{113} (2019), 3015--3030.
\url{https://doi.org/10.1007/s13398-019-00667-8}

\bibitem{BarbeiroFajardo}
A.~S.~V.~Barbeiro and R.~A.~dos Santos Fajardo,
\emph{A construction via forcing of a hereditarily weakly Koszmider space},
Rev. Roumaine Math. Pures Appl. \textbf{69} (2024), no.~1, 61--76.
\url{https://doi.org/10.59277/RRMPA.2024.61.76}

\bibitem{Engelking}
R.~Engelking,
\emph{General Topology}, 2nd ed., Sigma Series in Pure Mathematics
\textbf{6}, Heldermann, Berlin, 1989.
\url{https://www.heldermann.de/SSPM/SSPM06/sspm06.htm}

\bibitem{Fedorchuk}
V.~V.~Fedorchuk,
\emph{Fully closed mappings and the consistency of some theorems of general topology with the axioms of set theory},
Math. USSR-Sb. \textbf{28} (1976), no.~1, 1--26.
\url{https://doi.org/10.1070/SM1976v028n01ABEH001637}

\bibitem{Glodkowski}
D.~G{\l}odkowski,
\emph{A Banach space $C(K)$ reading the dimension of $K$},
J. Funct. Anal. \textbf{285} (2023), no.~4, 109986.
\url{https://doi.org/10.1016/j.jfa.2023.109986}

\bibitem{Hart}
K.~P.~Hart,
\emph{Efimov's problem}, in \emph{Open Problems in Topology II},
E.~Pearl (ed.), Elsevier, Amsterdam, 2007, 171--177.
\url{https://research.tudelft.nl/en/publications/efimovs-problem/}

\bibitem{Jarchow}
H.~Jarchow,
\emph{Locally Convex Spaces}, Teubner, Stuttgart, 1981.
\url{https://doi.org/10.1007/978-3-322-90559-8}

\bibitem{Jensen}
R.~B.~Jensen,
\emph{The fine structure of the constructible hierarchy},
Ann. Math. Logic \textbf{4} (1972), 229--308.
\url{https://doi.org/10.1016/0003-4843(72)90001-0}

\bibitem{Josefson}
B.~Josefson,
\emph{Weak sequential convergence in the dual of a Banach space does not imply norm convergence},
Ark. Mat. \textbf{13} (1975), 79--89.
\url{https://doi.org/10.1007/BF02386198}

\bibitem{KKLPS}
J.~K\k{a}kol, W.~Kubi\'s, M.~L\'opez-Pellicer and D.~Sobota,
\emph{Descriptive Topology in Selected Topics of Functional Analysis:
Updated and Expanded}, 2nd ed., Developments in Mathematics \textbf{24},
Springer, Cham, 2025.
\url{https://doi.org/10.1007/978-3-031-76062-4}

\bibitem{KSZ}
J.~K\k{a}kol, D.~Sobota and L.~Zdomskyy,
\emph{Grothendieck $C(K)$-spaces and the Josefson--Nissenzweig theorem},
Fund. Math. \textbf{263} (2023), no.~2, 105--131.
\url{https://doi.org/10.4064/fm218-6-2023}
Author version (with the cited theorem numbering):
\url{https://arxiv.org/abs/2207.13990}

\bibitem{KS2018}
J.~K\k{a}kol and W.~\'Sliwa,
\emph{Efimov spaces and the separable quotient problem for spaces $C_p(K)$},
J. Math. Anal. Appl. \textbf{457} (2018), 104--113.
\url{https://doi.org/10.1016/j.jmaa.2017.08.010}

\bibitem{KaniaKakol2DCP}
T.~Kania and J.~K\k{a}kol,
\emph{Talagrand compacta, 2DCP, and pointwise quotients},
revised manuscript, 2026.

\bibitem{Koszmider}
P.~Koszmider,
\emph{Banach spaces of continuous functions with few operators},
Math. Ann. \textbf{330} (2004), no.~1, 151--183.
\url{https://doi.org/10.1007/s00208-004-0544-z}

\bibitem{KoszmiderSilber}
P.~Koszmider and Z.~Silber,
\emph{Countably tight dual ball with a nonseparable measure},
J. Lond. Math. Soc. \textbf{110} (2024), no.~4, e12988.
\url{https://doi.org/10.1112/jlms.12988}

\bibitem{Lacey}
H.~E.~Lacey,
\emph{Separable quotients of Banach spaces},
An. Acad. Brasil. Ci\^enc. \textbf{44} (1972), 185--189.

\bibitem{Nissenzweig}
A.~Nissenzweig,
\emph{$w^*$ sequential convergence},
Israel J. Math. \textbf{22} (1975), 266--272.
\url{https://doi.org/10.1007/BF02761594}

\bibitem{Rinot}
A.~Rinot,
\emph{Jensen's diamond principle and its relatives},
in \emph{Set Theory and Its Applications},
Contemp. Math. \textbf{533}, American Mathematical Society, Providence, RI, 2011, 125--156.
\url{https://doi.org/10.48550/arXiv.0911.2151}

\bibitem{Rosenthal}
H.~P.~Rosenthal,
\emph{On quasi-complemented subspaces of Banach spaces, with an appendix on compactness of operators from $L^p(\mu)$ to $L^r(\nu)$},
J. Funct. Anal. \textbf{4} (1969), 176--214.
\url{https://doi.org/10.1016/0022-1236(69)90011-1}

\bibitem{Talagrand}
M.~Talagrand,
\emph{Un nouveau $C(K)$ qui poss\`ede la propri\'{e}t\'{e} de Grothendieck},
Israel J. Math. \textbf{37} (1980), 181--191.
\url{https://doi.org/10.1007/BF02762879}

\end{thebibliography}
\end{document}
